\documentclass[12pt, letterpaper]{article}
\usepackage[margin=1in]{geometry}
\usepackage[T1]{fontenc}
\usepackage{lmodern}
\usepackage{cmap}
\usepackage{amsmath, amssymb, amsthm}
\usepackage{mathtools}
\usepackage{enumitem}
\usepackage{booktabs}
\usepackage{xcolor}
\usepackage{titlesec}
\usepackage{parskip}
\usepackage{hyperref}
\usepackage{fancyhdr}
\usepackage{tcolorbox}
\usepackage{graphicx}
\tcbuselibrary{skins, breakable}

\definecolor{darkblue}{RGB}{15,55,115}
\definecolor{midblue}{RGB}{30,100,180}
\definecolor{theorybg}{RGB}{245,250,255}
\definecolor{questionbg}{RGB}{255,252,240}
\definecolor{evencolor}{RGB}{60,0,100}

\hypersetup{colorlinks=true,linkcolor=darkblue,urlcolor=midblue,
  pdftitle={Kurlat --- A Course in Modern Macroeconomics --- Complete Solutions, Chapters 1--4}}

\setlength{\headheight}{14pt}
\pagestyle{fancy}
\fancyhf{}
\renewcommand{\headrulewidth}{0.4pt}
\fancyhead[L]{\small\textcolor{darkblue}{\textit{Kurlat --- A Course in Modern Macroeconomics}}}
\fancyhead[R]{\small\textcolor{darkblue}{\thepage}}
\fancyfoot[C]{\small\textcolor{gray}{Complete Solutions --- Chapters 1--4}}

\titleformat{\section}[block]
  {\large\bfseries\color{darkblue}}{}{0pt}{}
  [\vspace{2pt}\color{darkblue}\rule{\textwidth}{1.5pt}\vspace{4pt}]

\titleformat{\subsection}[block]
  {\normalsize\bfseries\color{midblue}}{}{0pt}{}
  [\vspace{1pt}\color{midblue}\rule{0.5\textwidth}{0.8pt}\vspace{3pt}]

\tcbset{
  theorybox/.style={
    enhanced,breakable,colback=theorybg,colframe=darkblue,
    fonttitle=\bfseries\small\color{darkblue},boxrule=0.5pt,arc=3pt,
    left=6pt,right=6pt,top=4pt,bottom=4pt,title={Key Theory},
  },
  questionbox/.style={
    enhanced,breakable,colback=questionbg,colframe=gray!60,
    fonttitle=\bfseries\small\color{gray!60!black},boxrule=0.4pt,arc=2pt,
    left=6pt,right=6pt,top=4pt,bottom=4pt,title={\textit{Question}},
  },
  oddbox/.style={
    enhanced,breakable,colback=white,colframe=darkblue,
    attach boxed title to top left={yshift=-2mm,xshift=6mm},
    boxed title style={colback=darkblue,colframe=darkblue,arc=2pt,
      boxrule=0pt,left=6pt,right=6pt,top=3pt,bottom=3pt},
    fonttitle=\bfseries\large\color{white},
    boxrule=0.8pt,arc=3pt,
    left=8pt,right=8pt,top=10pt,bottom=6pt,
  },
  evenbox/.style={
    enhanced,breakable,colback=white,colframe=evencolor,
    attach boxed title to top left={yshift=-2mm,xshift=6mm},
    boxed title style={colback=evencolor,colframe=evencolor,arc=2pt,
      boxrule=0pt,left=6pt,right=6pt,top=3pt,bottom=3pt},
    fonttitle=\bfseries\large\color{white},
    boxrule=0.8pt,arc=3pt,
    left=8pt,right=8pt,top=10pt,bottom=6pt,
  },
  databox/.style={
    enhanced,breakable,colback=gray!8,colframe=gray!55,
    fonttitle=\bfseries\small\color{gray!60!black},boxrule=0.4pt,arc=2pt,
    left=6pt,right=6pt,top=4pt,bottom=4pt,title={Data exercise --- no closed-form answer},
  },
}

\newcommand{\pd}[2]{\dfrac{\partial #1}{\partial #2}}
\newcommand{\dd}[2]{\dfrac{d #1}{d #2}}
\newcommand{\MRS}{\mathrm{MRS}}
\newcommand{\MU}{\mathrm{MU}}
\newcommand{\Var}{\mathrm{Var}}
\newcommand{\Cov}{\mathrm{Cov}}
\newcommand{\E}{\mathrm{E}}
\newcommand{\MPK}{\mathrm{MPK}}
\newcommand{\MPL}{\mathrm{MPL}}
\newcommand{\stst}{^{\ast}}

\setlist[enumerate,1]{label=(\alph*),itemsep=6pt,parsep=2pt}
\setlist[enumerate,2]{label=(\roman*),itemsep=3pt}

\begin{document}

%======================================================================
% TITLE PAGE
%======================================================================
\begin{titlepage}
\centering
\vspace*{2cm}
{\Huge\bfseries\color{darkblue} A Course in Modern Macroeconomics\par}
\vspace{0.6cm}
{\Large\color{midblue} Pablo Kurlat (2020)\par}
\vspace{1.5cm}
{\LARGE\bfseries Complete Solutions\par}
\vspace{0.3cm}
{\Large Chapters 1--4 \quad ($23$ exercises)\par}
\vspace{2cm}

\begin{tcolorbox}[theorybox,title={Scope and course mapping},width=0.88\textwidth]
\small
Macroeconomia I --- MPE Insper, 2026/T3 --- Prof.\ Pedro G.\ Duarte.\\[4pt]
These four chapters cover \textbf{Aula 1} (Kurlat ch.\ 1--2) and \textbf{Aula 2}
(Kurlat ch.\ 3 and \S4.1--4.2), plus \S4.3--4.5 which belong to \textbf{Aula 3}.
Together they are the source material for \textbf{Lista de Exerc\'icios 1}.
\end{tcolorbox}

\vspace{0.8cm}
\begin{tabular}{clcc}
\toprule
\textbf{Ch.} & \textbf{Title} & \textbf{Exercises} & \textbf{Pages} \\
\midrule
1 & GDP                              & 1.1--1.5   & 27--29 \\
2 & Beyond GDP                       & 2.1--2.10  & 40--44 \\
3 & Basic Facts about Economic Growth & 3.1--3.3   & 51--52 \\
4 & The Solow Growth Model           & 4.1--4.5   & 72--74 \\
\midrule
\multicolumn{2}{r}{\textbf{Total}} & \textbf{23} & \\
\bottomrule
\end{tabular}

\vfill
\begin{tcolorbox}[colback=red!4,colframe=red!50,arc=2pt,width=0.88\textwidth,
  title={\bfseries\small On verification --- please read},fonttitle=\color{red!60!black}]
\small
\textbf{Kurlat provides no answer key.} The book has no ``Answers to Odd-Numbered
Problems'' section and no published solutions manual; it ends with the bibliography
(pp.~318--319). Every result below was therefore verified by other means, and the
method used is stated explicitly at each step:
\begin{itemize}[leftmargin=1.2em,itemsep=2pt,topsep=3pt]
\item \textbf{Executed numerical check} --- marked \checkmark. Every number was computed
      in Python and the code is reproduced so you can re-run it.
\item \textbf{Internal consistency} --- the three GDP methods are made to agree; steady
      states are substituted back; limiting and special cases are checked.
\item \textbf{Cross-book} --- for canonical models (Solow, factor prices) the method is
      checked against Romer (2012) ch.~1 and its solutions manual, both in
      \texttt{Livros/}. Notation differs; conclusions must not.
\end{itemize}
Where an exercise requires downloading data (2.1, 3.1--3.3), no numerical answer is
asserted. Those are marked as \emph{data exercises}: the method, the expected
qualitative pattern and runnable code are given instead.
\end{tcolorbox}
\vspace{1cm}
\end{titlepage}

\tableofcontents
\newpage

%======================================================================
\section{Chapter 1 \quad GDP}
%======================================================================

\begin{tcolorbox}[theorybox,title={Chapter 1 Overview}]
\textbf{The three methods.} GDP is the market value of final goods and services produced
domestically in a period. It can be measured three equivalent ways:
\begin{itemize}[leftmargin=1.3em,itemsep=2pt,topsep=2pt]
\item \textbf{Production:} sum of \emph{value added} $=$ output $-$ intermediate inputs.
\item \textbf{Expenditure:} $Y = C + I + G + (X - M)$.
\item \textbf{Income:} wages $+$ profits $+$ interest $+$ rents $+$ indirect taxes $-$ subsidies.
\end{itemize}
The three must give the same number; forcing them to agree is the best check on any
accounting answer.

\medskip
\textbf{Traps.} Intermediate goods are netted out (no double counting). Inventory changes
are investment. Imports are subtracted because they appear inside $C$, $I$ or $G$ but were
not produced domestically. Transfers (taxes, pensions) are \emph{not} production. Government
output is valued at input cost.

\medskip
\textbf{Real vs.\ nominal.} Real GDP holds prices fixed. Using base-year prices
(a Laspeyres-type index) and current-year prices (Paasche-type) give \emph{different}
growth rates whenever relative prices move --- the substitution bias. The
\textbf{chain-weighted} measure takes the geometric (Fisher) average of adjacent-year
links, which is why national accounts use it.

\medskip
\textbf{Across countries.} Converting at market exchange rates understates poor countries'
output because non-tradables are cheap there. \textbf{PPP} values every country's basket at
a common set of prices.
\end{tcolorbox}

%---- Topic: GDP Accounting --- the Three Methods ----
\subsection{GDP Accounting --- the Three Methods}

%------- 1.1 -------
% Topic: GDP Accounting --- the Three Methods
% Keywords: value added, intermediate goods, inventories, imports, transfers, government output at cost
\begin{tcolorbox}[oddbox,title={Problem 1.1 \quad Accounting}]
\begin{tcolorbox}[questionbox]
How does GDP accounting record the following events? For each of them, describe how they
would be computed in GDP accounts using the income method, the production method and the
expenditure method.\\[3pt]
(a) A car manufacturer buys components from Japan for \$1 to be used in production later on
and stores them at its warehouse.\\
(b) A car manufacturer buys components from Japan for \$1 and uses half of those components
in the production of a car, that it sells to Andy for \$2. It stores the rest of the
components.\\
(c) An army battalion is deployed to the border to repel a threatened Canadian invasion.
The soldiers earn wages of \$10,000 and use ammunition that the government buys for \$5,000.
The ammunition is produced using \$2,000 of imported steel and 100 hours of work, for which
the workers were paid \$1,000.\\
(d) Walmart sells 1000 bottles of Coca-Cola for \$1,500. It had previously paid \$1,200 for
them.\\
(e) A shipyard builds a cruise ship. It pays wages of \$200,000, interest on loans (from US
residents) of \$100,000 and \$300,000 for imported raw materials. The ship is sold for
\$1,000,000 to a cruise company. In the same year, the cruise company has revenue for
\$50,000 from operating cruises, pays wages of \$20,000 to its workers and has no other
expenses. Half the cruise revenue comes from tourists who reside in the United States and
half comes from tourists who reside abroad.\\
(f) The government collects \$1000 in income taxes from Roger.\\
(g) Roger earns \$4000 for working as a babysitter and pays \$1000 in income taxes.
\end{tcolorbox}

\textbf{Solution:}

\emph{Strategy.} For each event compute all three ways and require them to agree. That
agreement is the verification --- there is no answer key to check against.

\begin{enumerate}

\item \textbf{Contribution to GDP: \$0.}

\begin{center}
\begin{tabular}{lll}
\toprule
\textbf{Method} & \textbf{Entries} & \textbf{Total} \\
\midrule
Production  & no domestic output; imported components are not US production & \$0 \\
Expenditure & $I$ (inventories) $= +\$1$; \quad $M = \$1 \Rightarrow -\$1$ & \$0 \\
Income      & no wages, profits or interest generated domestically & \$0 \\
\bottomrule
\end{tabular}
\end{center}

The purchase raises \emph{inventory investment} by \$1 and imports by \$1. The two cancel
exactly. This is the general principle: buying a foreign good never changes GDP at the
moment of purchase, whatever the buyer does with it.

\item \textbf{Contribution to GDP: \$1.50.}

Of the \$1 of components, \$0.50 is used up in production and \$0.50 is stored.

\begin{center}
\begin{tabular}{lll}
\toprule
\textbf{Method} & \textbf{Entries} & \textbf{Total} \\
\midrule
Production  & output \$2 $-$ intermediate inputs used \$0.50 & \$1.50 \\
Expenditure & $C = \$2$; \; $I = +\$0.50$ (unused components); \; $M = \$1$ & \$1.50 \\
Income      & profits $= 2 - 0.50 = \$1.50$ (no wages stated) & \$1.50 \\
\bottomrule
\end{tabular}
\end{center}

\checkmark\ $2 + 0.50 - 1 = 1.50$. Note that only the \emph{half actually consumed in
production} is netted out as an intermediate input. The stored half is still an asset of
the firm, so it sits in inventories.

\item \textbf{Contribution to GDP: \$13,000.}

Two producers here: the government (which produces defence services) and the ammunition
manufacturer.

\begin{center}
\begin{tabular}{lll}
\toprule
\textbf{Method} & \textbf{Entries} & \textbf{Total} \\
\midrule
Production  & government value added $=\$10{,}000$ (valued at input cost) & \\
            & ammunition maker $= 5{,}000 - 2{,}000 = \$3{,}000$ & \$13,000 \\
\midrule
Expenditure & $G = 10{,}000 + 5{,}000 = \$15{,}000$; \quad $M = \$2{,}000$ & \$13,000 \\
\midrule
Income      & soldiers' wages \$10,000 $+$ ammo workers' wages \$1,000 & \\
            & $+$ ammo firm profits $(5{,}000-2{,}000-1{,}000)=\$2{,}000$ & \$13,000 \\
\bottomrule
\end{tabular}
\end{center}

\checkmark\ Three methods agree. Two points worth stating explicitly:
\begin{itemize}[leftmargin=1.3em,itemsep=2pt,topsep=2pt]
\item \textbf{Government output is valued at input cost} because there is no market price
      for repelling an invasion. So the soldiers' \$10,000 of wages \emph{is} the measured
      value of the defence service produced.
\item The ammunition is an \textbf{intermediate input} into defence from the production
      side, but from the expenditure side it is a \emph{final} purchase by government. Both
      routes give \$13,000 --- they just split the same total differently.
\end{itemize}

\item \textbf{Contribution to GDP this year: \$300.}

\begin{center}
\begin{tabular}{lll}
\toprule
\textbf{Method} & \textbf{Entries} & \textbf{Total} \\
\midrule
Production  & retail value added $= 1{,}500 - 1{,}200 = \$300$ & \$300 \\
Expenditure & $C = \$1{,}500$; \; $I$ (inventories) $= -\$1{,}200$ & \$300 \\
Income      & Walmart's wages $+$ profits $= \$300$ & \$300 \\
\bottomrule
\end{tabular}
\end{center}

Walmart produces a \emph{retail service} --- getting the bottles onto a shelf near the
customer --- and that service is worth \$300. The \$1,200 of Coca-Cola was counted as GDP
when it was \emph{produced}; selling it now is a drawdown of inventories, which is exactly
why the expenditure side subtracts \$1,200.

\item \textbf{Contribution to GDP: \$750,000.}

\emph{Shipyard.} Value added $= 1{,}000{,}000 - 300{,}000 = \$700{,}000$.
Income check: wages $200{,}000 +$ interest $100{,}000 +$ profits
$(1{,}000{,}000-300{,}000-200{,}000-100{,}000) = 400{,}000$, total $\$700{,}000$. \checkmark

\emph{Cruise company.} No intermediate inputs, so value added $=$ revenue $=\$50{,}000$.
Income: wages $20{,}000 +$ profits $30{,}000 = \$50{,}000$. \checkmark

\begin{center}
\begin{tabular}{lll}
\toprule
\textbf{Method} & \textbf{Entries} & \textbf{Total} \\
\midrule
Production  & $700{,}000 + 50{,}000$ & \$750,000 \\
Expenditure & $I = \$1{,}000{,}000$ (the ship is a capital good) & \\
            & $C = \$25{,}000$ (US tourists); \; $X = \$25{,}000$ (foreign tourists) & \\
            & $M = \$300{,}000$ & \$750,000 \\
Income      & $700{,}000 + 50{,}000$ & \$750,000 \\
\bottomrule
\end{tabular}
\end{center}

\checkmark\ $1{,}000{,}000 + 25{,}000 + 25{,}000 - 300{,}000 = 750{,}000$.

Three things to notice: the ship is \textbf{investment}, not consumption, because it is a
durable asset bought by a firm to produce with. \textbf{Interest paid to US residents is
domestic income}, so it stays inside value added --- it is a distribution of the \$700,000,
not a cost that reduces it. And cruises sold to foreign residents are \textbf{exports} even
though the service is consumed inside the country.

\item \textbf{Contribution to GDP: \$0.}

A tax is a \textbf{transfer}: purchasing power moves from Roger to the government, but
nothing is produced. All three methods record zero. In the income method in particular,
taxes are not a new source of income --- income is measured \emph{before} the transfer, and
counting it again at the government would double count.

\item \textbf{Contribution to GDP: \$4,000.}

\begin{center}
\begin{tabular}{lll}
\toprule
\textbf{Method} & \textbf{Entries} & \textbf{Total} \\
\midrule
Production  & babysitting services produced, no intermediate inputs & \$4,000 \\
Expenditure & $C = \$4{,}000$ by the households who hired him & \$4,000 \\
Income      & Roger's labour income, \emph{pre-tax} & \$4,000 \\
\bottomrule
\end{tabular}
\end{center}

The \$1,000 of tax does not reduce GDP --- by (f) it is a transfer. Contrast this with the
\emph{same} babysitting done by a parent for their own child, which is real production but
is \textbf{excluded} from GDP because no market transaction records it. That boundary
problem is the subject of Problem 1.4.

\end{enumerate}

\begin{tcolorbox}[theorybox,title={Takeaway}]
Every line above obeys one rule: \textbf{GDP counts production, and only production, once}.
Imports are subtracted because they were produced elsewhere; intermediates are netted out
because they are already inside the final good; transfers are ignored because they move
claims, not output.
\end{tcolorbox}
\end{tcolorbox}

%---- Topic: Real GDP, Index Numbers and Chain Weighting ----
\subsection{Real GDP, Index Numbers and Chain Weighting}

%------- 1.2 -------
% Topic: Real GDP, Index Numbers and Chain Weighting
% Keywords: nominal GDP, real GDP, base year, Laspeyres, Paasche, chain weighting, PPP, market exchange rate
\begin{tcolorbox}[evenbox,title={Problem 1.2 \quad Comparisons Across Time and Across Countries}]
\begin{tcolorbox}[questionbox]
Suppose these are the prices (in US dollars) and quantities of goods A and B produced in the
US in 2017 and 2018:
\begin{center}
\begin{tabular}{ccccc}
\toprule
 & $p_A$ & $q_A$ & $p_B$ & $q_B$ \\
\midrule
2017 & 4 & 5  & 3 & 3 \\
2018 & 1 & 10 & 4 & 2 \\
\bottomrule
\end{tabular}
\end{center}
(a) What was nominal GDP in 2017?\\
(b) What was nominal GDP in 2018?\\
(c) What was real GDP in 2018 at 2017 prices (computed using fixed 2017 prices)? Using this
measure, how much did GDP grow between 2017 and 2018?\\
(d) What was real GDP in 2017 at 2018 prices (computed using fixed 2018 prices)? Using this
measure, how much did GDP grow between 2017 and 2018? What explains the difference between
the two measures?\\
(e) How much did GDP grow between 2017 and 2018 using the chain-weighted method?\\[4pt]
In Thailand, prices (in Thai baht) and quantities in 2017 were:
\begin{center}
\begin{tabular}{ccccc}
\toprule
 & $p_A$ & $q_A$ & $p_B$ & $q_B$ \\
\midrule
2017 & 30 & 1 & 10 & 2 \\
\bottomrule
\end{tabular}
\end{center}
(f) What was nominal GDP in Thailand in 2017, expressed in baht.\\
(g) Suppose the exchange rate in 2017 was 25 baht per dollar. What was GDP in Thailand in
2017, expressed in US dollars at market exchange rates?\\
(h) What was GDP in Thailand in 2017 at PPP? What accounts for the difference between the
market exchange rate measure and the PPP measure?\\
(i) What was the PPP exchange rate between the baht and the dollar?
\end{tcolorbox}

\textbf{Solution:}

\begin{enumerate}

\item $\text{GDP}^{nom}_{2017} = p_A^{17}q_A^{17} + p_B^{17}q_B^{17}
= 4\cdot 5 + 3\cdot 3 = 20 + 9 = \boxed{29}$ \checkmark

\item $\text{GDP}^{nom}_{2018} = 1\cdot 10 + 4\cdot 2 = 10 + 8 = \boxed{18}$ \checkmark

Nominal GDP \emph{fell} by $38\%$. Hold that thought --- it is about to be contradicted by
both real measures.

\item Value 2018 quantities at 2017 prices:
$$\text{GDP}^{real}_{2018}\big|_{p^{17}} = 4\cdot 10 + 3\cdot 2 = 40 + 6 = \boxed{46}$$
Growth relative to 2017 nominal (which is also 2017 at 2017 prices, $=29$):
$$g_L = \frac{46}{29} - 1 = 0.5862 = \boxed{58.62\%}$$ \checkmark

\item Value 2017 quantities at 2018 prices:
$$\text{GDP}^{real}_{2017}\big|_{p^{18}} = 1\cdot 5 + 4\cdot 3 = 5 + 12 = \boxed{17}$$
Growth, comparing to 2018 nominal ($=18$, i.e.\ 2018 at 2018 prices):
$$g_P = \frac{18}{17} - 1 = 0.0588 = \boxed{5.88\%}$$ \checkmark

\textbf{Why the two differ so much (58.6\% vs 5.9\%).} Good A collapsed in relative price
(from $p_A/p_B = 4/3$ to $1/4$) and simultaneously \emph{doubled} in quantity. That is the
textbook substitution pattern.
\begin{itemize}[leftmargin=1.3em,itemsep=2pt,topsep=2pt]
\item Using \textbf{2017 prices} gives good A a weight of 4 --- its old, high price. The good
      that grew fastest is weighted by the price it had \emph{before} it got cheap, so
      growth looks large. This is the \textbf{Laspeyres} (base-weighted) direction and it
      \textbf{overstates} growth.
\item Using \textbf{2018 prices} weights A at 1 --- its new, low price. The fast-growing good
      barely counts, so growth looks small. This is the \textbf{Paasche} direction and it
      \textbf{understates} growth.
\end{itemize}
Neither is ``right''. The gap between them is a genuine index-number problem, not an error.

\item The chain-weighted (Fisher) growth rate is the geometric mean of the two links:
$$1 + g_C = \sqrt{(1+g_L)(1+g_P)} = \sqrt{1.5862 \times 1.0588} = 1.2960$$
$$g_C = \boxed{29.60\%}$$ \checkmark

It sits between the two, as it must, and it is symmetric: it gives no privileged status to
either year's prices. This is why the BEA (and the IBGE) publish chained series.

\item $\text{GDP}^{Thai}_{2017} = 30\cdot 1 + 10\cdot 2 = 30 + 20 = \boxed{50 \text{ baht}}$ \checkmark

\item At the market rate of 25 baht per dollar:
$$\frac{50 \text{ baht}}{25 \text{ baht/\$}} = \boxed{\$2}$$ \checkmark

\item At PPP, value \emph{Thai quantities} at \emph{US prices} (the common price vector):
$$\text{GDP}^{Thai,PPP} = p_A^{US}q_A^{Thai} + p_B^{US}q_B^{Thai} = 4\cdot 1 + 3\cdot 2
= 4 + 6 = \boxed{\$10}$$ \checkmark

\textbf{The PPP measure is five times the market-exchange-rate measure.} The reason: at the
market rate, Thai goods are being valued at \emph{Thai} prices converted to dollars, and Thai
prices are low. Good A costs 30 baht $= \$1.20$ in Thailand but \$4 in the US; good B costs
10 baht $= \$0.40$ versus \$3. Thailand's output is physically the same whatever prices we
use to value it, and PPP prices it consistently with US output. Market exchange rates are
driven by the supply and demand for \emph{tradables} and for financial assets, so they
systematically understate the real output of poor countries, where non-tradables (housing,
haircuts, domestic services) are cheap. This is the Balassa--Samuelson logic.

\item The PPP exchange rate is the rate that equalises the cost of the same basket:
$$e^{PPP} = \frac{\text{GDP in baht}}{\text{GDP in PPP dollars}} = \frac{50}{10}
= \boxed{5 \text{ baht per dollar}}$$ \checkmark

Compare with the market rate of 25 baht/\$. The baht is ``undervalued'' by a factor of 5 on
this measure --- equivalently, the price level in Thailand is one fifth of the US level.
\end{enumerate}

\begin{tcolorbox}[theorybox,title={Verification code}]
\footnotesize
\begin{verbatim}
pA17,qA17,pB17,qB17 = 4,5,3,3
pA18,qA18,pB18,qB18 = 1,10,4,2
n17 = pA17*qA17 + pB17*qB17          # 29
n18 = pA18*qA18 + pB18*qB18          # 18
r18 = pA17*qA18 + pB17*qB18          # 46  -> gL = 46/29-1 = 58.62%
r17 = pA18*qA17 + pB18*qB17          # 17  -> gP = 18/17-1 =  5.88%
gC  = ((r18/n17)*(n18/r17))**0.5 - 1 # 29.60%
th  = 30*1 + 10*2                    # 50 baht
ppp = pA17*1 + pB17*2                # $10 ; market: 50/25 = $2
\end{verbatim}
\end{tcolorbox}
\end{tcolorbox}

%------- 1.3 -------
% Topic: Real GDP, Index Numbers and Chain Weighting
% Keywords: chain weighting, Fisher index, substitution bias, base year, average growth rate
\begin{tcolorbox}[oddbox,title={Problem 1.3 \quad Chained GDP}]
\begin{tcolorbox}[questionbox]
The country of Fructus produces Apples, Bananas and Cherries. Its production statistics are
given in the book's table for 2000--2020 (also downloadable as a spreadsheet from the book
website).\\[3pt]
(a) Compute a real GDP series at year-2000 prices using base year prices and using the
chained method.\\
(b) Plot both series and comment on any differences.\\
(c) What was the average growth rate according to each method?
\end{tcolorbox}

\textbf{Solution:}

\begin{enumerate}

\item Two constructions, both normalised to nominal GDP in 2000 ($=9{,}500$):

\emph{Fixed base-year prices.} $\;Y^{base}_t = \sum_i p_{i,2000}\, q_{i,t}$.

\emph{Chained (Fisher).} Build year-on-year links and cumulate:
$$L_t = \frac{\sum_i p_{i,t-1}q_{i,t}}{\sum_i p_{i,t-1}q_{i,t-1}}, \qquad
P_t = \frac{\sum_i p_{i,t}q_{i,t}}{\sum_i p_{i,t}q_{i,t-1}}, \qquad
Y^{chain}_t = Y^{chain}_{t-1}\sqrt{L_t P_t}$$

Selected values (full series from the code below):

\begin{center}
\begin{tabular}{crrr}
\toprule
\textbf{Year} & \textbf{Nominal} & \textbf{Base-2000} & \textbf{Chained} \\
\midrule
2000 &  9,500.0 &  9,500.0 &  9,500.0 \\
2001 &  9,782.0 &  9,650.0 &  9,640.0 \\
2005 & 10,021.0 & 10,400.0 & 10,341.8 \\
2010 & 10,812.0 & 11,230.0 & 10,950.0 \\
2015 & 12,048.0 & 12,390.0 & 11,922.8 \\
2020 & 13,127.0 & 13,960.0 & 12,976.2 \\
\bottomrule
\end{tabular}
\end{center}

\item \textbf{The two series separate steadily, with the base-year series always above.}
By 2020 the fixed-base measure reads $13{,}960$ against $12{,}976$ for the chained one --- a
gap of about $7.6\%$ that opens up gradually rather than jumping.

The mechanism is exactly the one flagged in Problem 1.2(d). Look at what happens to
cherries: their price falls from 20 to 12 (a $40\%$ decline) while their quantity rises from
200 to 361 (a $+81\%$ increase). Bananas do the opposite --- price up from 30 to 61,
quantity down from 50 to 34. Holding 2000 prices fixed keeps valuing cherries at their old,
high price of 20 forever, so the fastest-growing good keeps an inflated weight and measured
growth is too high. The chained index re-weights every year and tracks the moving relative
prices.

\item Average annual growth over the 20 year span, $g = (Y_{2020}/Y_{2000})^{1/20} - 1$:

\begin{center}
\begin{tabular}{lrr}
\toprule
\textbf{Method} & \textbf{2020 level} & \textbf{Average growth} \\
\midrule
Base-2000 prices & 13,960.0 & $\mathbf{1.9432\%}$ \\
Chained          & 12,976.2 & $\mathbf{1.5713\%}$ \\
\midrule
Difference       &          & $0.3718$ p.p.\ per year \\
\bottomrule
\end{tabular}
\end{center}
\checkmark\ Both computed in Python; code below.

\textbf{Interpretation.} A gap of $0.37$ percentage points a year sounds small and is not:
compounded over the two decades it is the entire $7.6\%$ level discrepancy. This is the
substitution bias in its purest form, and it is precisely why the US switched from
fixed-base to chained real GDP in 1996 --- the old method was overstating growth in exactly
the sectors (computers, electronics) where prices were falling fastest and quantities rising
fastest.
\end{enumerate}

\begin{tcolorbox}[theorybox,title={Verification code (executed)}]
\footnotesize
\begin{verbatim}
import numpy as np
q = np.array([qA, qB, qC], float).T      # (21, 3) quantities, 2000..2020
p = np.array([pA, pB, pC], float).T      # (21, 3) prices

base = (q * p[0]).sum(axis=1)            # fixed 2000 prices

chain = np.empty(len(q)); chain[0] = (q[0]*p[0]).sum()
for t in range(1, len(q)):
    L = (q[t]*p[t-1]).sum() / (q[t-1]*p[t-1]).sum()   # Laspeyres link
    P = (q[t]*p[t]).sum()   / (q[t-1]*p[t]).sum()     # Paasche link
    chain[t] = chain[t-1] * np.sqrt(L*P)              # Fisher link

n = len(q) - 1
print((base[-1]/base[0])**(1/n) - 1)     # 0.019432
print((chain[-1]/chain[0])**(1/n) - 1)   # 0.015713
\end{verbatim}
\end{tcolorbox}
\end{tcolorbox}

%------- 1.5 -------
% Topic: Real GDP, Index Numbers and Chain Weighting
% Keywords: relative prices, substitution bias, technological progress, Baumol cost disease, demand elasticity
\begin{tcolorbox}[oddbox,title={Problem 1.5 \quad Changes in Relative Prices}]
\begin{tcolorbox}[questionbox]
We saw that using an earlier year as the base year to compute real GDP results in a high
rate of growth if the sectors that are expanding most are those whose relative price is
falling. Can you think of reasons why that should be the case (i.e., economic forces that
make the same types of goods become relatively cheaper and be produced in higher
quantities)? Can you think of reasons why the opposite should be the case (i.e., economic
forces that make the same types of goods become relatively expensive and be produced in
higher quantities)?
\end{tcolorbox}

\textbf{Solution:}

The question is really about whether the observed price-quantity pair is generated by a
\textbf{supply shift} or a \textbf{demand shift}. Sliding along a stable demand curve gives
falling price with rising quantity; sliding along a stable supply curve gives rising price
with rising quantity.

\begin{enumerate}

\item \textbf{Cheaper \emph{and} more of it --- supply-driven (the common case).}

\begin{itemize}[leftmargin=1.3em,itemsep=3pt,topsep=2pt]
\item \textbf{Sector-biased technological progress.} Productivity growth is wildly uneven
      across sectors. Semiconductors, computers, telecoms and increasingly solar panels have
      had enormous productivity gains, which push the supply curve out: unit cost and
      therefore relative price fall, and quantity rises. This is the single most important
      reason and the one that motivated the switch to chained indices.
\item \textbf{Learning by doing and economies of scale.} Cumulative production lowers unit
      cost, which lowers price, which raises quantity, which lowers cost further --- a
      self-reinforcing loop.
\item \textbf{Elastic demand.} For the quantity response to be \emph{large}, demand must be
      price-elastic. Goods with close substitutes and discretionary timing (consumer
      electronics) qualify; necessities do not.
\item \textbf{Trade liberalisation and offshoring} shift the effective supply curve out for
      importables.
\end{itemize}

\item \textbf{More expensive \emph{and} more of it --- demand-driven.}

\begin{itemize}[leftmargin=1.3em,itemsep=3pt,topsep=2pt]
\item \textbf{Income-elastic goods (luxuries).} As income grows, demand for healthcare,
      education, higher-end housing and personal services shifts out faster than supply. Both
      price and quantity rise.
\item \textbf{Baumol's cost disease.} In sectors where productivity growth is inherently
      slow --- a string quartet needs the same four players it needed in 1800 --- wages
      still have to track the economy-wide average or workers leave. Unit costs, and hence
      relative prices, rise. If demand is income-elastic (health, education, live
      performance), quantity rises \emph{too}. This is the cleanest example of the
      combination the question asks about.
\item \textbf{Binding supply constraints.} Urban land, waterfront property, rare earths: the
      supply curve is steep, so a demand shift moves price much more than quantity --- but
      quantity does still rise.
\item \textbf{Network effects / quality ladders.} A good becomes more valuable as more people
      use it, shifting demand out while the price rises.
\end{itemize}

\end{enumerate}

\textbf{The index-number consequence.} Base-year weighting overstates growth in case (1) and
\emph{understates} it in case (2), because there the expanding sectors are the ones whose
prices are rising, so old prices under-weight them. Since case (1) empirically dominates in
modern economies --- technology-intensive goods are both the fastest-growing and the
fastest-cheapening --- fixed-base real GDP has an upward bias overall. Problem 1.3 is a
constructed illustration of exactly this.

\medskip
\emph{Connection to the course:} the same asymmetry reappears in
\texttt{rules/07\_moeda\_inflacao.md} as the \textbf{substitution bias in the CPI/IPCA} ---
a Laspeyres consumer price index overstates inflation for the same reason a Laspeyres
quantity index overstates growth.
\end{tcolorbox}

%---- Topic: The Boundaries of GDP ----
\subsection{The Boundaries of GDP}

%------- 1.4 -------
% Topic: The Boundaries of GDP
% Keywords: production boundary, illegal activities, imputation, comparability, national accounts conventions
\begin{tcolorbox}[evenbox,title={Problem 1.4 \quad Drugs and Prostitution}]
\begin{tcolorbox}[questionbox]
Read the following article:
\url{http://www.economist.com/news/finance-and-economics/21603073-italys-inclusion-illicit-activities-its-figures-excites-much-interest-sex}.
What do you think? Should drug production and prostitution be included in the calculation of
GDP?
\end{tcolorbox}

\textbf{Solution:}

\begin{tcolorbox}[databox,title={Note on sourcing}]
This exercise sends you to a specific 2014 \emph{Economist} piece, which I have not
retrieved. The background it refers to is the EU's ESA~2010 accounting standard, which from
September 2014 required member states to include illegal activities --- drugs, prostitution,
smuggling --- in GDP where they are voluntary market transactions. Italy's measured GDP rose
by roughly one percent as a result. The argument below is the economics; check the
article for its specific framing before quoting it.
\end{tcolorbox}

\textbf{Position: yes, they should be included} --- and the reasons are about consistency,
not about morality.

\begin{enumerate}

\item \textbf{The production boundary is defined by transactions, not by approval.} GDP
already counts tobacco, alcohol, gambling, weapons and payday lending. It counts legal
prostitution in the Netherlands and Germany and legal cannabis in Colorado. If the criterion
were social desirability, the accounts would need a moral weighting for every line item, and
no two countries would agree on it. The criterion actually used --- was there a voluntary
exchange producing something someone paid for --- is at least well defined.

\item \textbf{Comparability collapses without it.} This is the strongest argument. Cannabis
is legal in some EU countries and not others; prostitution is legal in Germany and illegal
in France. If illegal activity is excluded, then Germany's GDP includes its sex industry and
France's does not, and the two numbers are not measuring the same thing. Worse, a country
that \emph{legalises} an activity would show a spurious jump in GDP with no change in real
activity. Uruguay's cannabis legalisation and Portugal's decriminalisation would both create
artefacts. Including illegal activity everywhere removes the discontinuity.

\item \textbf{Coherence of the accounts.} The inputs to illegal industries are often already
counted --- the chemicals, the rent, the imported precursors. Counting inputs but not output
puts a hole in the value-added identity.

\item \textbf{Policy relevance.} Debt-to-GDP ratios, EU budget contributions and deficit
limits are all denominated in GDP. Italy's revision changed its measured debt ratio without
changing a euro of debt. That is an argument for measuring the denominator consistently, not
for leaving parts of it out.

\end{enumerate}

\textbf{The serious objections, and what they actually establish:}

\begin{enumerate}[label=(\roman*)]
\item \textbf{Measurement is guesswork.} There are no invoices. Estimates are built from
      seizure data, survey-based prevalence and assumed prices, and revisions are large.
      \emph{But} --- this is an argument about \emph{precision}, not about \emph{principle}.
      Owner-occupied housing rent is also imputed rather than observed, and it is a far
      larger share of GDP than drugs.
\item \textbf{Coercion breaks the ``voluntary'' premise.} Trafficking and coerced
      prostitution are not voluntary exchange, and treating them as consensual market output
      is a genuine category error. The honest response is that the accounts cannot separate
      the coerced from the voluntary component, which is a real and unresolved defect.
\item \textbf{The welfare reading is wrong.} Counting drug sales as output while ignoring
      addiction, overdose and crime overstates wellbeing. This is correct --- but it is not
      special to illegal goods. GDP does not net out pollution, commuting or the harm from
      any legal product either.
\end{enumerate}

\textbf{Conclusion.} Objection (iii) is the deepest, and it does not argue for excluding
illegal activity from GDP. It argues that \textbf{GDP is not a welfare measure and should
not be read as one} --- which is exactly the thesis of Chapter 2 and of the Jones--Klenow
$\lambda$. Include the activity in the production measure; use a separate welfare measure
when the question is whether people are better off.
\end{tcolorbox}

%======================================================================
\section{Chapter 2 \quad Beyond GDP}
%======================================================================

\begin{tcolorbox}[theorybox,title={Chapter 2 Overview}]
\textbf{The question.} GDP per capita measures production. Welfare depends on consumption,
how it is distributed, how long people live and how much they work. Chapter 2 builds
measures that price those things.

\medskip
\textbf{The HDI} is a simple geometric average of three normalised sub-indices (income,
education, life expectancy). It is transparent but the weights are arbitrary --- there is no
theory saying a year of life expectancy is worth a given amount of income.

\medskip
\textbf{The Jones--Klenow $\lambda$} fixes exactly that by grounding the weights in
preferences. Behind a Rawlsian veil of ignorance, ask: by what factor $\lambda$ would we have
to scale \emph{US} consumption to make a person indifferent between being dropped at random
into the US and into country $j$? That $\lambda$ is an \textbf{equivalent variation} and is
the country's welfare relative to the US. The book's eq.\ (2.2.1) is
$$u(c,l,a) = \E\left[\bar u + \frac{c^{1-\sigma}}{1-\sigma} - \theta(1-l)^2\right] a$$
and $\lambda$ solves $\;u(\lambda c_{US}, l_{US}, a_{US}) = u(c_j, l_j, a_j)$.

\medskip
\textbf{The role of $\sigma$.} The curvature parameter does double duty. It is risk aversion
behind the veil, and because a concave $u$ penalises spread, it is simultaneously the
\emph{inequality aversion} of the welfare measure. Higher $\sigma$ means a more unequal
country is marked down harder. As $\sigma \to \infty$ the criterion converges on Rawls's
difference principle: only the worst-off person counts.

\medskip
\textbf{The empirical punchline.} Western Europe looks \emph{better} on $\lambda$ than on
GDP per capita (more leisure, longer lives, less inequality); much of Africa looks
\emph{worse}, driven mainly by low life expectancy.
\end{tcolorbox}

%---- Topic: The HDI and the Jones--Klenow Welfare Measure ----
\subsection{The HDI and the Jones--Klenow Welfare Measure}

%------- 2.1 -------
% Topic: The HDI and the Jones--Klenow Welfare Measure
% Keywords: HDI, equivalent variation, relative GDP, scatterplot, welfare-to-GDP ratio
\begin{tcolorbox}[oddbox,title={Problem 2.1 \quad Comparing the HDI and the Jones \& Klenow Welfare Measure}]
\begin{tcolorbox}[questionbox]
Download the UN data that goes into building the HDI from
\url{http://hdr.undp.org/en/data} and the Jones \& Klenow Data from
\url{http://web.stanford.edu/~chadj/papers.html\#rawls}. For each country, construct:
(a) Relative GDP: $\dfrac{\text{GDP per capita}}{\text{GDP per capita in the US}}$;
(b) Welfare-to-GDP: $\dfrac{\lambda}{\text{Relative GDP}}$;
(c) HDI-to-GDP: $\dfrac{\text{HDI}}{\text{Relative GDP}}$;
and plot a scatterplot of Welfare-to-GDP against HDI-to-GDP. What do each of these ratios
measure? Is the impression we get from the Jones \& Klenow measure very different from the
one we get from the HDI? What countries look better under each measure and why?
\end{tcolorbox}

\textbf{Solution:}

\begin{tcolorbox}[databox]
This requires two external datasets. No numbers are asserted below --- the method, the
expected pattern (which follows from the structure of the two measures and from the book's
own Figure 2.2.3) and runnable code are given instead.
\end{tcolorbox}

\begin{enumerate}

\item \textbf{What each ratio measures.}

\begin{itemize}[leftmargin=1.3em,itemsep=3pt,topsep=2pt]
\item \textbf{Relative GDP} is the pure income yardstick, normalised so the US $=1$.
\item \textbf{Welfare-to-GDP} $=\lambda/(\text{relative GDP})$ asks: \emph{once we account
for consumption share, inequality, leisure and mortality, does this country do better
or worse than its income alone would suggest?} A ratio above 1 means the country
punches above its income; below 1 means income flatters it.
\item \textbf{HDI-to-GDP} asks the same question using the UN's weights instead of
preference-based ones.
\end{itemize}

Both are \textbf{residual} measures --- they strip out income and show what the non-income
components add. That is what makes them comparable on a scatterplot.

\item \textbf{Do the two give the same impression? Partly --- and the disagreements are
systematic.}

Expect a \textbf{positive but loose correlation}. Both load heavily on life expectancy, so
countries with long lives relative to income (much of Latin America, southern Europe) score
above 1 on both, and countries with an HIV/AIDS-driven mortality deficit (South Africa,
Botswana) score below 1 on both. That shared component drives the positive slope.

The disagreements come from what each measure includes that the other does not:

\begin{center}
\begin{tabular}{p{0.30\textwidth}p{0.28\textwidth}p{0.28\textwidth}}
\toprule
\textbf{Component} & \textbf{In $\lambda$?} & \textbf{In HDI?} \\
\midrule
Life expectancy   & yes, priced by preferences & yes, arbitrary weight \\
Consumption level & yes (consumption, not GDP) & indirectly, via income \\
\textbf{Inequality} & \textbf{yes}, via $\sigma$ & \textbf{no} \\
\textbf{Leisure / hours} & \textbf{yes}, via $\theta$ & \textbf{no} \\
Education & no & yes, one third of the index \\
\bottomrule
\end{tabular}
\end{center}

Therefore:
\begin{itemize}[leftmargin=1.3em,itemsep=3pt,topsep=2pt]
\item \textbf{Better under $\lambda$ than under HDI:} \emph{Western Europe}. Short working
hours and low inequality are rewarded by $\lambda$ and completely invisible to the
HDI. This is Jones and Klenow's headline finding --- France and Italy look much closer
to the US on welfare than on income.
\item \textbf{Better under HDI than under $\lambda$:} countries with \textbf{strong schooling
relative to consumption and high inequality} --- several post-socialist and some Latin
American economies. The HDI credits their educational attainment; $\lambda$ ignores
education entirely and penalises their income dispersion.
\item \textbf{Worse under $\lambda$:} high-inequality middle-income countries generally.
\textbf{Brazil} and \textbf{South Africa} are the clearest cases: the HDI's income
component uses the \emph{mean}, while $\lambda$'s concave $u$ effectively uses
something closer to a certainty-equivalent, which is far below the mean when dispersion
is large.
\end{itemize}

\item \textbf{The methodological point the exercise is driving at.} The HDI's weights are
chosen by committee; $\lambda$'s are derived from a stated utility function and can be
argued with on their own terms --- you can disagree about $\sigma$, $\theta$ or $\bar u$ and
recompute. That transparency is the real advantage, more than any particular ranking.
Conversely $\lambda$ omits education, which most people would say matters, so neither
dominates.

\end{enumerate}

\begin{tcolorbox}[theorybox,title={Code to run once you have the two files}]
\footnotesize
\begin{verbatim}
import pandas as pd, matplotlib.pyplot as plt

hdi = pd.read_csv("hdi.csv")            # country, hdi
jk  = pd.read_csv("jones_klenow.csv")   # country, lambda, gdp_pc

df = hdi.merge(jk, on="country").dropna()
us = df.loc[df.country == "United States", "gdp_pc"].iloc[0]

df["rel_gdp"]       = df.gdp_pc / us
df["welfare_ratio"] = df["lambda"] / df.rel_gdp
df["hdi_ratio"]     = df.hdi       / df.rel_gdp

ax = df.plot.scatter("hdi_ratio", "welfare_ratio", alpha=.6)
ax.axhline(1, ls=":", c="gray"); ax.axvline(1, ls=":", c="gray")
print(df[["country","rel_gdp","welfare_ratio","hdi_ratio"]]
        .sort_values("welfare_ratio", ascending=False).head(15))
print(df[["welfare_ratio","hdi_ratio"]].corr())
\end{verbatim}
\end{tcolorbox}
\end{tcolorbox}

%------- 2.2 -------
% Topic: The HDI and the Jones--Klenow Welfare Measure
% Keywords: omitted variables in welfare, revealed preference, compensating differentials, willingness to pay
\begin{tcolorbox}[evenbox,title={Problem 2.2 \quad Other Things that Matter}]
\begin{tcolorbox}[questionbox]
The welfare measure $\lambda$ takes into account data on consumption levels, inequality,
leisure and mortality.\\[3pt]
(a) Name one other variable that might be an important determinant of welfare that is not
included in standard GDP calculations.\\
(b) What data would you need in order to figure out how much weight to give to this
variable? Describe how you would use that data to come up with the correct way to include
the variable in question in the welfare calculation. What choices that people make might
reveal how much they care about this variable? Don't worry too much about the feasibility of
the data-collection procedure, but think carefully about why the observed choices would be
informative.
\end{tcolorbox}

\textbf{Solution:}

\begin{tcolorbox}[theorybox,title={This is question 2 of Lista de Exerc\'icios 1 (2 pts)}]
The Lista assigns this exercise directly. The grader will be looking for the
\textbf{revealed-preference logic} in part (b), not for a list of nice-sounding variables in
part (a). Spend your words on (b).
\end{tcolorbox}

\begin{enumerate}

\item \textbf{Candidate: personal security --- the risk of being a victim of violent crime.}

It is a good choice for Brazil specifically: homicide rates vary enormously across states and
municipalities while incomes do not vary nearly as much, which is what makes the weight
estimable. Other defensible answers: air and water quality, commuting time, political
freedom and civil liberties, job security, or the quality of social relationships. Any of
these is fine \emph{provided} you can describe in (b) how a choice would reveal its value.

Note that GDP treats crime perversely: spending on locks, private security, alarms and
insurance \emph{adds} to GDP, so a more dangerous country can look richer.

\item \textbf{How to price it: find a choice where people trade the variable against
money.} Stated preference (``how much would you pay to be safe?'') is unreliable because
nothing is at stake. Revealed preference uses actual decisions with real consequences.

\textbf{Method 1 --- hedonic pricing of housing (the workhorse).}

Two otherwise identical apartments, one in a neighbourhood with a higher homicide rate, will
not command the same rent. The rent gap capitalises the value of the safety difference.
Estimate
$$\ln(\text{rent}_i) = \beta_0 + \beta_1\,\text{crime}_i + \gamma' X_i + \varepsilon_i$$
where $X_i$ controls for size, age, amenities, school quality and distance to work. Then
$-\beta_1$ is the marginal willingness to pay for a unit reduction in crime risk, as a share
of housing cost.

\emph{Why this is informative:} a household choosing where to live is making a real,
costly, repeated trade-off between money and safety. It has every incentive to get it right.

\emph{Identification problem, and the fix:} crime is correlated with everything else that
makes a neighbourhood bad, so $\beta_1$ picks up more than safety. You need variation in
crime that is unrelated to other neighbourhood characteristics --- for example a policing
reform rolled out in some districts and not others, or the timing of a UPP-style
intervention. Compare rents \emph{within} the same neighbourhood before and after.

\textbf{Method 2 --- compensating wage differentials.}

Dangerous occupations pay more, all else equal. Regressing wages on fatality risk across
occupations, controlling for skill, gives the wage premium per unit of mortality risk and
hence a \textbf{value of a statistical life}. This is the same logic Jones and Klenow use to
price life expectancy, which is why it slots into their framework naturally.

\textbf{Method 3 --- migration.}

People leaving a violent region for a lower-paying job elsewhere reveal a lower bound on
what they will pay for safety.

\item \textbf{How to put it into the welfare calculation.} Extend the flow utility with a
security argument $s$:
$$u(c, l, a, s) = \left[\bar u + \frac{c^{1-\sigma}}{1-\sigma} - \theta(1-l)^2
+ \psi\,v(s)\right] a$$
and calibrate $\psi$ so that the model's marginal rate of substitution matches the estimated
one:
$$\left.\frac{\partial u/\partial s}{\partial u/\partial c}\right|_{\text{model}}
\;=\; \underbrace{\text{WTP per unit of safety}}_{\text{from the hedonic regression}}$$
Then recompute $\lambda$ with the extra term. Countries safer than the US get $\lambda$
revised up; Brazil, with homicide rates many times the US level, would be revised
\textbf{down}.

\emph{A caveat worth stating:} if crime mainly \emph{kills} people, its cost is already
partly captured through life expectancy $a$, and adding $\psi v(s)$ on top would double
count. The security term should be calibrated to capture fear, avoidance behaviour and
non-fatal victimisation --- the part that does \emph{not} show up in mortality.

\end{enumerate}
\end{tcolorbox}

%------- 2.9 -------
% Topic: The HDI and the Jones--Klenow Welfare Measure
% Keywords: risk aversion, inequality aversion, comparative statics of sigma, Western Europe
\begin{tcolorbox}[oddbox,title={Problem 2.9 \quad Risk Aversion and Welfare in Western Europe}]
\begin{tcolorbox}[questionbox]
If instead $\sigma = 1$, Jones and Klenow had set a higher value of $\sigma$, what would this
do to measured welfare in Western European countries? Explain.
\end{tcolorbox}

\textbf{Solution:}

\textbf{Answer: measured welfare in Western Europe would rise further relative to the US.}

The reasoning has two steps.

\begin{enumerate}

\item \textbf{$\sigma$ is simultaneously risk aversion and inequality aversion.} Behind the
veil of ignorance, an individual faces a lottery over who they will be. A more concave $u$
(higher $\sigma$) makes them dislike that lottery's spread more. Equivalently, the
certainty-equivalent consumption
$$c^{CE} = \left((1-\sigma)\,\E\!\left[\tfrac{c^{1-\sigma}}{1-\sigma}\right]\right)^{\frac{1}{1-\sigma}}$$
falls further below the mean as $\sigma$ rises, so unequal countries are penalised more
heavily.

\item \textbf{Western Europe is markedly less unequal than the US.} Gini coefficients for
disposable income run roughly $0.25$--$0.32$ in France, Germany, Belgium and the Nordics
against about $0.39$ in the US. So raising $\sigma$ imposes a \emph{larger} penalty on the
US than on Europe. Since $\lambda$ is measured \emph{relative to the US}, and the
denominator is being marked down more than the numerator, European $\lambda$ rises.

\end{enumerate}

\textbf{Reinforcing channel.} Higher $\sigma$ also raises the weight on the $\bar u$ term
relative to the consumption term --- that is, on \emph{being alive} versus consuming.
Western Europe's life expectancy exceeds the US, so this pushes the same way.

\textbf{Second reinforcing channel.} Europe's consumption \emph{level} is below the US. With
higher $\sigma$ the utility function flattens at high $c$, so differences in consumption
levels matter less at the margin and Europe's income shortfall is discounted --- again
favouring Europe. All three channels point the same way, which is why the conclusion is
robust.

\medskip
\textbf{Limiting case as a check.} As $\sigma \to \infty$ the measure converges on the
Rawlsian difference principle: welfare is determined by the worst-off person alone
(Problem 2.10 makes this precise). On that criterion Western European welfare states, with
their higher consumption floor, would rank \emph{well above} the US. The comparative static
is just the continuous version of that limit. \checkmark
\end{tcolorbox}

%---- Topic: Valuing Life Expectancy and the Composition of Consumption ----
\subsection{Valuing Life Expectancy and the Composition of Consumption}

%------- 2.3 -------
% Topic: Valuing Life Expectancy and the Composition of Consumption
% Keywords: NIPA, healthcare share, intermediate vs final consumption, double counting in welfare
\begin{tcolorbox}[oddbox,title={Problem 2.3 \quad Healthcare}]
\begin{tcolorbox}[questionbox]
When we calculate consumption, one of the (many) categories of consumption is healthcare
services.\\[3pt]
(a) Look up how much healthcare is consumed in the United States per year for a recent year.
State the total dollar amount and also what fraction of GDP and what fraction of consumption
is accounted for by healthcare. A good place to look for this data is the National Income and
Product Accounts (NIPA) at \url{https://apps.bea.gov/iTable/index_nipa.cfm}.\\
(b) Suppose that we are calculating welfare in the style of Jones and Klenow, taking into
account the impact of life expectancy on utility. Should we therefore subtract consumption of
healthcare services from our measure of consumption? What do you think?
\end{tcolorbox}

\textbf{Solution:}

\begin{enumerate}

\item \begin{tcolorbox}[databox,title={Look this up yourself --- NIPA Table 2.3.5}]
Navigate to NIPA \textbf{Table 2.3.5}, ``Personal Consumption Expenditures by Major Type of
Product'', and read off the \emph{Health care} line under Services. Divide by GDP
(Table 1.1.5, line 1) and by total PCE.

As an order of magnitude to sanity-check what you find: US health spending has been running
around $17$--$18\%$ of GDP on the broadest (national health expenditure) definition, and
health care is roughly a fifth of personal consumption expenditure. The PCE health-care line
is \emph{narrower} than total national health expenditure --- it excludes investment in
medical structures and equipment and some public health activity --- so expect a smaller
number from Table 2.3.5 than the headline ``18\% of GDP'' figure you will see quoted. Report
whichever you use and say which definition it is.
\end{tcolorbox}

\item \textbf{No --- do not subtract it. But the question is sharper than it looks, and the
right answer is ``it depends what health care buys''.}

\textbf{The case for subtracting (the ``intermediate good'' view).} Health care is arguably
not something people want for its own sake. Nobody enjoys chemotherapy. It is an
\emph{input} into producing health and longevity. If the Jones--Klenow measure already
rewards a country for its life expectancy $a$, then also counting the consumption that
\emph{bought} that life expectancy counts the same benefit twice --- once as the output
(longer life) and once as the input (the spending). On this view health care is like the
intermediate goods netted out in Chapter 1.

\textbf{Why that argument fails, or at least overreaches:}

\begin{itemize}[leftmargin=1.3em,itemsep=3pt,topsep=2pt]
\item \textbf{It proves far too much.} By the same logic you would subtract food (an input
into not starving), housing (an input into shelter) and exercise. Essentially all
consumption is an input into something. The consumption measure would collapse.
\item \textbf{Health care delivers more than longevity.} A hip replacement, cataract surgery,
insulin, dental work and pain management mostly improve the \emph{quality} of life,
not its length. That benefit appears \emph{nowhere} in $u(c,l,a)$ --- not in $a$, and
not in $c$ if you strip health care out. Subtracting would make the measure worse, not
better.
\item \textbf{The double counting is not symmetric across countries.} The US spends far more
per unit of life expectancy than other rich countries. If anything this is an argument
that US consumption is \emph{overstated} as a welfare measure --- i.e.\ that a dollar
of US health spending buys less wellbeing than a dollar elsewhere --- which is a
statement about \emph{prices and efficiency}, not about whether the category belongs
in the accounts.
\end{itemize}

\textbf{The defensible position.} Keep health care in consumption, and recognise that the
genuine problem is a \textbf{price-measurement} one, not a boundary one. If US health care is
expensive rather than abundant, then deflating US health consumption by an appropriate price
index --- one that accounts for quality and outcomes --- would reduce measured real US
consumption. That is where the correction belongs. Netting the whole category out is a blunt
instrument that throws away the quality-of-life benefits along with the double count.

\end{enumerate}
\end{tcolorbox}

%------- 2.4 -------
% Topic: Valuing Life Expectancy and the Composition of Consumption
% Keywords: consumption vs GDP, composition of output, defence spending, welfare accounting
\begin{tcolorbox}[evenbox,title={Problem 2.4 \quad Construction and the Armed Forces}]
\begin{tcolorbox}[questionbox]
Suppose we compared welfare in two otherwise identical countries. Country A dedicates 50\% of
GDP to maintaining a large army. Country B dedicates 50\% of GDP to building houses. In which
of the two countries would welfare be higher, according to the measure used by Jones and
Klenow. What do you think of this?
\end{tcolorbox}

\textbf{Solution:}

\begin{enumerate}

\item \textbf{What the measure says: welfare is the same in both, and low in both.}

This is the key mechanical point. The Jones--Klenow $\lambda$ is built on
\textbf{consumption}, not GDP. Both countries devote $50\%$ of output to something that is
not current consumption:
\begin{itemize}[leftmargin=1.3em,itemsep=2pt,topsep=2pt]
\item Country A: $G = 0.5Y$ on defence, so $C = 0.5Y$.
\item Country B: $I = 0.5Y$ on housing, so $C = 0.5Y$.
\end{itemize}
Since the two countries are ``otherwise identical'' --- same leisure, same life expectancy,
same inequality --- and consumption is identical, $\lambda_A = \lambda_B$. Both are measured
as \emph{worse off} than an otherwise identical country that consumed its whole output.

\item \textbf{What I think of this: the measure is right about A versus B, and wrong about
the level --- but for different reasons in each case, and the housing case is the more
serious defect.}

\textbf{On defence (country A), the measure is defensible.} An army is not consumed and
yields no direct utility. It may be \emph{necessary} --- protecting the country from
conquest --- but that is precisely the point: it is a cost of maintaining the ability to
consume, not a benefit. A country that must spend half its output on defence genuinely is
worse off than an identical country that faces no threat. Treating it as forgone consumption
captures something real. This is the same logic as the ``regrettables'' discussion in the
welfare literature: commuting, prisons and pollution abatement all absorb resources without
generating utility.

\textbf{On housing (country B), the measure is clearly deficient.} Houses are durable capital
that yields a \textbf{flow of housing services for decades}. In a single-period snapshot
country B looks identical to country A, but this is a timing artefact:
\begin{itemize}[leftmargin=1.3em,itemsep=3pt,topsep=2pt]
\item In the \textbf{next} period, B's residents enjoy a larger stock of houses, and national
accounts \emph{do} count that as consumption via \textbf{imputed rent} on
owner-occupied housing. So over a full horizon B is unambiguously better off.
\item The defect is therefore \textbf{the static, one-period framing}, not the use of
consumption. Investment today is consumption tomorrow. Comparing two countries by
current consumption alone penalises the one that is building its future --- the same
objection one would make to ranking a country as ``worse off'' for having a high saving
rate.
\end{itemize}

\item \textbf{The general lesson.} Any welfare measure built on a single period's consumption
mis-ranks economies that differ in their \emph{investment} rates. The right fix is to
evaluate the \textbf{present discounted value} of the consumption stream, or equivalently to
compare steady states rather than transition points --- which is exactly what the Solow model
of Chapter 4 forces you to do, and what the Golden Rule (\S4.3) is designed to answer.

\end{enumerate}
\end{tcolorbox}

%------- 2.5 -------
% Topic: Valuing Life Expectancy and the Composition of Consumption
% Keywords: revealed preference, value of statistical life, hedonic pricing, identification, confounding
\begin{tcolorbox}[oddbox,title={Problem 2.5 \quad Measuring the Value of Life Expectancy}]
\begin{tcolorbox}[questionbox]
Would each of the following observations be useful in determining how much to weigh life
expectancy in measured welfare? Explain.\\[3pt]
(a) Observing how life expectancy varies across people with different income levels within a
country.\\
(b) Observing how life expectancy varies across countries.\\
(c) Observing how much people are willing to pay for funerals.\\
(d) Observing the difference in house prices in neighborhoods with different murder rates.
\end{tcolorbox}

\textbf{Solution:}

\emph{The criterion throughout:} to price life expectancy we need a situation in which
someone \textbf{chooses}, at a real cost, between money and mortality risk. Correlations
that arise from something \emph{other} than such a choice do not identify the weight.

\begin{enumerate}

\item \textbf{Partly useful, but badly confounded --- causality runs both ways.}

Rich people do live longer. But the correlation cannot be read as ``people buy longevity with
income'' because:
\begin{itemize}[leftmargin=1.3em,itemsep=2pt,topsep=2pt]
\item \textbf{Reverse causality:} healthier people earn more. Illness reduces
labour supply and productivity, so part of the gradient runs from health to income.
\item \textbf{Common causes:} education, discount rates and family background drive both.
A patient, educated person both saves more and smokes less.
\end{itemize}
There is a further, subtler problem: even a \emph{causal} income--longevity gradient measures
how much longevity people \emph{get}, not how much they \emph{value} it, unless we know they
were deliberately trading one for the other.

\emph{Verdict:} suggestive, not identifying. Would need an income shock unrelated to health
--- a lottery win, an inheritance, a pension reform --- to be usable.

\item \textbf{Least useful of the four.}

Cross-country variation in life expectancy is driven by public health systems, sanitation,
epidemiology, conflict and climate --- essentially none of which is an individual's
willingness-to-pay decision. Nobody \emph{chose} to be born in a country with a high
under-five mortality rate. Confounding is at its worst here because rich countries differ
from poor ones in every dimension at once.

\emph{Verdict:} useful for \emph{describing} the differences that the welfare measure needs
to weigh, but useless for determining the \emph{weight} itself. It is the data you plug in,
not the data you calibrate with.

\item \textbf{Not useful, and it is a category error.}

Funeral spending is consumption by \textbf{the survivors}, not by the deceased. It measures
willingness to pay for a ritual, for social status and for grief management. It says nothing
about how much a living person values an additional year of their own life --- the quantity
$a$ in the utility function. Funeral spending would still be positive in a world where nobody
valued longevity at all.

\emph{Verdict:} no. This is the distractor in the list, and identifying it as such is the
point of including it.

\item \textbf{The most useful of the four --- this is the right kind of data.}

House prices in high- versus low-murder-rate neighbourhoods reflect a \textbf{real,
voluntary, costly choice} in which households trade money directly against mortality risk.
That is exactly the trade-off we need to price. The method is hedonic regression:
$$\ln(\text{price}_i) = \beta_0 + \beta_1\,\text{murder rate}_i + \gamma' X_i + \varepsilon_i$$
Converting the price differential into a per-unit-of-risk figure yields a \textbf{value of a
statistical life}, which is precisely what calibrates $\bar u$ in the Jones--Klenow utility
function.

\emph{The identification caveat, which you should state:} murder rates correlate with school
quality, amenities, employment access and neighbourhood income, so a naive $\beta_1$ is
biased. Credible estimates exploit sharp variation --- policing interventions, boundary
discontinuities, or before/after comparisons within the same area.

\emph{Related and equally valid:} compensating wage differentials across occupations with
different fatality rates. Same logic, different market.

\end{enumerate}

\begin{tcolorbox}[theorybox,title={Ranking}]
$(d) \;>\; (a) \;>\; (b) \;>\; (c)$. The ordering tracks one thing: how close the observation
is to an actual choice in which a person gave up money for a lower chance of dying.
\end{tcolorbox}
\end{tcolorbox}

%------- 2.6 -------
% Topic: Valuing Life Expectancy and the Composition of Consumption
% Keywords: equivalent variation, lambda, value of a life-year, isoelastic utility, Jones-Klenow
\begin{tcolorbox}[evenbox,title={Problem 2.6 \quad The Value of Life Expectancy}]
\begin{tcolorbox}[questionbox]
In the US, average consumption per capita is about \$35,000. Life expectancy is 79 years, so
in the Jones and Klenow experiment, Rawls would have a 79\% chance of being alive. In
Bolivia, average consumption per capita is \$3,700 and life expectancy is 68 years. Assume
that there is no inequality in either country so that everyone who is alive gets the average
level of consumption. Assume the following utility function:
$$u(c,a) = \left[\bar u + \frac{c^{1-\sigma}}{1-\sigma}\right] a$$
with $\sigma = 0.5$ and $\bar u = 48.1$.\\[3pt]
(a) Suppose you made a US resident the following offer: you can buy a health plan that costs
$x$ dollars per year and will extend your lifetime by one year. What is the price $x$ that
would make them indifferent between getting the health plan or not?\\
(b) How much would a Bolivian resident be willing to pay for such a plan?\\
(c) Write down the special case of equation (2.2.1) that applies to this example and solve
for $\lambda$ for Bolivia. Don't replace any numbers yet, leave it in terms of $a_{US}$,
$a_{Bol}$, $c_{US}$, $c_{Bol}$, $\sigma$ and $\bar u$.\\
(d) Replace the values of $a_{US}$, $a_{Bol}$, $c_{US}$, $c_{Bol}$, $\sigma$ and $\bar u$ to
find a value for $\lambda$. How does it compare to $\dfrac{c_{Bol}}{c_{US}}$? Why?
\end{tcolorbox}

\textbf{Solution:}

\emph{Setup.} With $\sigma = 1/2$ the consumption term simplifies usefully:
$$\frac{c^{1-\sigma}}{1-\sigma} = \frac{c^{1/2}}{1/2} = 2\sqrt{c}$$
so the \textbf{flow} utility is $\;\phi(c) \equiv \bar u + 2\sqrt{c}\;$ and total utility is
$u(c,a) = \phi(c)\,a$.

\begin{center}
\begin{tabular}{lrrr}
\toprule
& $c$ & $a$ & $\phi(c) = 48.1 + 2\sqrt{c}$ \\
\midrule
US      & 35,000 & 79 & $422.2657$ \\
Bolivia &  3,700 & 68 & $169.7553$ \\
\bottomrule
\end{tabular}
\end{center}

\begin{enumerate}

\item \textbf{The US resident pays $x$ so that living 80 years on $c-x$ equals living 79
years on $c$:}
$$\left[\bar u + 2\sqrt{c_{US}-x}\right] \cdot 80 = \left[\bar u + 2\sqrt{c_{US}}\right]\cdot 79$$
Solve for $x$ in closed form:
$$x = c_{US} - \left(\frac{1}{2}\left[\tfrac{79}{80}\phi(c_{US}) - \bar u\right]\right)^{2}$$
Numerically:
$$\tfrac{79}{80}\cdot 422.2657 = 416.9874,\qquad
\tfrac{1}{2}(416.9874 - 48.1) = 184.4437,\qquad 184.4437^2 = 34{,}019.48$$
$$\boxed{x_{US} = 35{,}000 - 34{,}019.48 = \$980.52} \qquad \checkmark$$
That is about $2.80\%$ of annual consumption.

\item \textbf{Same calculation for Bolivia}, with $68 \to 69$ years:
$$\left[\bar u + 2\sqrt{3700-x}\right]\cdot 69 = 169.7553 \cdot 68$$
$$\tfrac{68}{69}\cdot 169.7553 = 167.2951,\qquad \tfrac{1}{2}(167.2951-48.1) = 59.5976,
\qquad 59.5976^2 = 3{,}551.86$$
$$\boxed{x_{Bol} = 3{,}700 - 3{,}551.86 = \$148.14} \qquad \checkmark$$

\textbf{Comparison and the economics.} In \emph{absolute} terms the American pays $6.6$ times
more --- unsurprising, since they are $9.5$ times richer. But as a \emph{share of
consumption}:
$$\frac{x_{US}}{c_{US}} = 2.80\%, \qquad \frac{x_{Bol}}{c_{Bol}} = 4.00\%$$
The Bolivian gives up a \emph{larger fraction} of consumption. Why? Because they start with
a shorter life. Going from 68 to 69 years is a proportionally bigger gain
($1/68 = 1.47\%$) than 79 to 80 ($1/79 = 1.27\%$), and because $u$ is concave in $c$ but
\emph{linear} in $a$, the proportional gain in years is what matters. The absolute gap,
meanwhile, reflects that a year of life is worth more when each year is spent consuming more.

\item \textbf{The special case of (2.2.1).} Drop leisure (not modelled here) and the
expectation over inequality (assumed away). Equation (2.2.1) becomes
$$u(c,a) = \left[\bar u + \frac{c^{1-\sigma}}{1-\sigma}\right]a$$
and $\lambda$ is defined by making Rawls indifferent between Bolivia and a rescaled US:
$$\boxed{\;a_{US}\left[\bar u + \frac{(\lambda c_{US})^{1-\sigma}}{1-\sigma}\right]
= a_{Bol}\left[\bar u + \frac{c_{Bol}^{\,1-\sigma}}{1-\sigma}\right]\;}$$
Solving for $\lambda$:
$$\frac{(\lambda c_{US})^{1-\sigma}}{1-\sigma}
= \frac{a_{Bol}}{a_{US}}\left[\bar u + \frac{c_{Bol}^{\,1-\sigma}}{1-\sigma}\right] - \bar u$$
$$\boxed{\;\lambda = \frac{1}{c_{US}}\left\{(1-\sigma)\left(\frac{a_{Bol}}{a_{US}}
\left[\bar u + \frac{c_{Bol}^{\,1-\sigma}}{1-\sigma}\right] - \bar u\right)\right\}^{\frac{1}{1-\sigma}}\;}$$

\item \textbf{Plugging in} $\sigma = 0.5$, $\bar u = 48.1$, $a_{US}=79$, $a_{Bol}=68$,
$c_{US}=35{,}000$, $c_{Bol}=3{,}700$:
$$\frac{a_{Bol}}{a_{US}} = \frac{68}{79} = 0.86076, \qquad \phi(c_{Bol}) = 169.7553$$
$$0.86076 \times 169.7553 - 48.1 = 146.1243 - 48.1 = 98.0243$$
$$\lambda c_{US} = \left(\frac{98.0243}{2}\right)^{2} = 49.0122^2 = 2{,}402.19$$
$$\boxed{\lambda = \frac{2{,}402.19}{35{,}000} = 0.0686} \qquad \checkmark$$

\textbf{Comparison with relative consumption:}
$$\frac{c_{Bol}}{c_{US}} = \frac{3{,}700}{35{,}000} = 0.1057
\qquad\text{versus}\qquad \lambda = 0.0686, \qquad
\frac{\lambda}{c_{Bol}/c_{US}} = 0.649$$

\textbf{Why $\lambda$ is well below relative consumption.} Bolivian welfare is only $6.9\%$
of US welfare even though Bolivian consumption is $10.6\%$ of US consumption --- roughly a
third lower. The entire wedge comes from \textbf{life expectancy}: Bolivians live 68 years
against 79, a factor of $0.861$, and in this utility function years multiply the flow, so the
mortality shortfall scales welfare down directly. Consumption alone \emph{understates} how
much worse off Bolivia is.

This is the central empirical message of Jones and Klenow and the reason $\lambda$ is worth
computing at all: for poor countries, and especially for those with an AIDS-driven or
conflict-driven mortality deficit, GDP per capita is not merely an imperfect welfare proxy
--- it is systematically \emph{too optimistic}.

\end{enumerate}

\begin{tcolorbox}[theorybox,title={Verification code (executed)}]
\footnotesize
\begin{verbatim}
from scipy.optimize import brentq
sig, ubar = 0.5, 48.1
flow = lambda c: ubar + c**(1-sig)/(1-sig)      # = ubar + 2*sqrt(c)
u    = lambda c, a: flow(c)*a
cUS, aUS, cBol, aBol = 35000, 79, 3700, 68

x_us  = brentq(lambda x: u(cUS-x, aUS+1) - u(cUS, aUS), 0, cUS-1)      #  980.52
x_bol = brentq(lambda x: u(cBol-x, aBol+1) - u(cBol, aBol), 0, cBol-1) #  148.14

lam = brentq(lambda l: aUS*flow(l*cUS) - aBol*flow(cBol), 1e-9, 1)     # 0.068626
# closed form agrees to 6 dp:
lam_cf = ((1-sig)/cUS**(1-sig) * ((aBol/aUS)*flow(cBol) - ubar))**(1/(1-sig))
\end{verbatim}
\end{tcolorbox}
\end{tcolorbox}

%---- Topic: Risk Aversion, Inequality and the Veil of Ignorance ----
\subsection{Risk Aversion, Inequality and the Veil of Ignorance}

%------- 2.7 -------
% Topic: Risk Aversion, Inequality and the Veil of Ignorance
% Keywords: CRRA, certainty equivalent, risk premium, calibration of sigma, indifference
\begin{tcolorbox}[oddbox,title={Problem 2.7 \quad Measuring Risk Aversion}]
\begin{tcolorbox}[questionbox]
The students of Concave University are all identical in their ability and preferences, which
are well described by the function $u(c) = \E\left[\dfrac{c^{1-\sigma}}{1-\sigma}\right]$.
Two employers recruit CU graduates. Stabilis offers each graduate a fixed salary of \$50,000.
Lotteris instead offers them a base salary of \$20,000 plus a bonus scheme that depends on
their performance, which is entirely driven by luck. The bonus is either \$10,000 or
\$100,000, with equal probability. The job itself is the same in both firms. About half the
CU graduates choose to work for Stabilis and half prefer Lotteris, and they all say they
found it hard to choose because both offers were similarly attractive. What value of $\sigma$
is consistent with their decisions?
\end{tcolorbox}

\textbf{Solution:}

\begin{enumerate}

\item \textbf{Translate the story into an indifference condition.} ``They found it hard to
choose because both offers were similarly attractive'' means the graduates are
\textbf{indifferent} --- that is the whole content of the setup, and it is what pins down
$\sigma$.

Lotteris pays base $+$ bonus, so the two outcomes are
$$20{,}000 + 10{,}000 = 30{,}000 \quad\text{or}\quad 20{,}000 + 100{,}000 = 120{,}000,
\qquad \text{each with probability } \tfrac12$$
Note the expected pay at Lotteris is $\tfrac12(30{,}000) + \tfrac12(120{,}000) = \$75{,}000$,
far above Stabilis's \$50,000. Lotteris must pay a \textbf{risk premium} of \$25,000 to
attract half the class --- and the size of that premium is exactly what reveals $\sigma$.

\item \textbf{Set up the equation.}
$$\frac{50{,}000^{\,1-\sigma}}{1-\sigma}
= \frac{1}{2}\cdot\frac{30{,}000^{\,1-\sigma}}{1-\sigma}
+ \frac{1}{2}\cdot\frac{120{,}000^{\,1-\sigma}}{1-\sigma}$$
The $\frac{1}{1-\sigma}$ factor cancels \emph{provided} $\sigma \ne 1$ --- but be careful
with the sign, since $1-\sigma < 0$ when $\sigma > 1$. Safer to keep it and solve
numerically. The condition says that \$50,000 is the \textbf{certainty equivalent} of the
Lotteris lottery.

\item \textbf{Solve.} There is no closed form; solve numerically. Bracketing:
\begin{center}
\begin{tabular}{lrrl}
\toprule
$\sigma$ & $u(50{,}000)$ & $\E[u_{Lotteris}]$ & Preferred \\
\midrule
$0.50$ & $447.214$ & $519.615$ & Lotteris \\
$1.00$ & $10.8198$ & $11.0021$ & Lotteris \\
$\mathbf{1.7964}$ & $-2.2732\!\times\!10^{-4}$ & $-2.2732\!\times\!10^{-4}$ & \textbf{indifferent} \\
$2.00$ & $-2.0000\!\times\!10^{-5}$ & $-2.0833\!\times\!10^{-5}$ & Stabilis \\
\bottomrule
\end{tabular}
\end{center}
$$\boxed{\sigma \approx 1.796} \qquad \checkmark$$

\item \textbf{Interpretation.} A $\sigma$ of about $1.8$ is squarely inside the range
macroeconomists normally use --- calibrations typically sit between 1 and 4, and Jones and
Klenow themselves use $\sigma = 1$. So the CU students' behaviour is entirely ordinary.

Two things worth noting:
\begin{itemize}[leftmargin=1.3em,itemsep=3pt,topsep=2pt]
\item \textbf{$\sigma$ is a curvature parameter, and its magnitude is only interpretable
relative to the size of the gamble.} This is a very large gamble --- a $4\times$ spread
between outcomes --- and it takes only moderate curvature to generate a \$25,000
premium on it.
\item \textbf{The monotonicity is the useful check.} At low $\sigma$ everyone takes the
lottery (its mean is 50\% higher); at high $\sigma$ everyone takes the safe job. There
is exactly one crossing, so $\sigma$ is identified. If the problem had said
\emph{everyone} chose Stabilis we would only get an inequality, $\sigma > 1.796$.
\end{itemize}

\item \textbf{Connection to Chapter 2's real purpose.} This same $\sigma$ is the inequality
aversion in the welfare measure. A society where consumption is $30{,}000$ for half the
population and $120{,}000$ for the other half is, on these preferences, exactly as good as
one where everyone has $50{,}000$ --- even though mean consumption in the first is $75{,}000$.
The \$25,000 gap \emph{is} the welfare cost of that inequality. Problem 2.8 pushes this
further.

\end{enumerate}

\begin{tcolorbox}[theorybox,title={Verification code (executed)}]
\footnotesize
\begin{verbatim}
from scipy.optimize import brentq
import numpy as np
safe, lo, hi = 50000, 20000+10000, 20000+100000     # 50000 ; 30000 ; 120000

def gap(s):
    if abs(s-1) < 1e-12:
        return np.log(safe) - 0.5*(np.log(lo) + np.log(hi))
    f = lambda c: c**(1-s)/(1-s)
    return f(safe) - 0.5*(f(lo) + f(hi))

print(brentq(gap, 0.05, 5))        # 1.796394
\end{verbatim}
\end{tcolorbox}
\end{tcolorbox}

%------- 2.8 -------
% Topic: Risk Aversion, Inequality and the Veil of Ignorance
% Keywords: inequality aversion, mean-preserving spread, expected utility, certainty equivalent
\begin{tcolorbox}[evenbox,title={Problem 2.8 \quad Inequality and Risk Aversion}]
\begin{tcolorbox}[questionbox]
Compare the following two countries. Both have a population of 100. Within each country, we
label individuals in order of increasing consumption. In country A, the consumption of
individual $j$ is $c_A(j) = 100 + 8j$, while in country B the consumption of individual $j$
is $c_B(j) = 200 + 4j$.\\[3pt]
(a) Plot the consumption patterns of each country, with an individual's label $j$ (which
ranges from 1 to 100) on the horizontal axis and their consumption on the vertical axis.\\
(b) Compute per capita consumption in each country. [Note: if you want, you can approximate
sums with integrals]\\
(c) Suppose the utility function in both countries is $u(c) = \dfrac{c^{1-\sigma}}{1-\sigma}$.
For what values of $\sigma$ is expected utility higher in country A? Interpret your results.
\end{tcolorbox}

\textbf{Solution:}

\begin{enumerate}

\item \textbf{Both are straight lines in $j$; A is steeper and starts lower.}

\begin{center}
\begin{tabular}{lrrrr}
\toprule
& intercept & slope & $c(1)$ & $c(100)$ \\
\midrule
Country A & 100 & 8 & 108 & 900 \\
Country B & 200 & 4 & 204 & 600 \\
\bottomrule
\end{tabular}
\end{center}

They cross where $100+8j = 200+4j \Rightarrow 4j = 100 \Rightarrow j = 25$, at
$c = 300$. So:
\begin{itemize}[leftmargin=1.3em,itemsep=2pt,topsep=2pt]
\item the \textbf{poorest quarter} ($j < 25$) is better off in \textbf{B};
\item the \textbf{richest three quarters} ($j > 25$) are better off in \textbf{A}.
\end{itemize}
This crossing is the whole exercise: neither country dominates, so the ranking must depend on
how much we care about the bottom --- i.e.\ on $\sigma$.

\item \textbf{Per capita consumption.} Using $\sum_{j=1}^{100} j = \frac{100\cdot 101}{2} = 5050$,
so $\bar j = 50.5$:
$$\bar c_A = 100 + 8(50.5) = 100 + 404 = \boxed{504}, \qquad
\bar c_B = 200 + 4(50.5) = 200 + 202 = \boxed{402}$$
\checkmark\ (The integral approximation $\int_0^{100}$ would give $\bar j = 50$ and hence
$500$ and $400$ --- close, and fine for the qualitative point.)

\textbf{Country A is $25.4\%$ richer on average, but far more unequal:} standard deviations
are $230.9$ for A against $115.5$ for B, exactly double, since the spread scales with the
slope. A ranking based on mean consumption alone would pick A without hesitation. The point of
the exercise is that this is not obviously right.

\item \textbf{Country A is preferred for $\sigma < \sigma^{\ast} \approx 2.072$, and B
above it.}

Compute $\E[u] = \frac{1}{100}\sum_{j=1}^{100} \frac{c(j)^{1-\sigma}}{1-\sigma}$ for each:

\begin{center}
\begin{tabular}{crrc}
\toprule
$\sigma$ & $\E[u_A]$ & $\E[u_B]$ & Preferred \\
\midrule
$0$    & $504$ & $402$ & A \\
$0.25$ & $138.715$ & $118.743$ & A \\
$0.50$ & $43.5329$ & $39.6651$ & A \\
$0.75$ & $18.4992$ & $17.7631$ & A \\
$1$ (log) & $6.08797$ & $5.95172$ & A \\
$1.5$  & $-0.0993398$ & $-0.10323$ & A \\
$2$    & $-2.70274\!\times\!10^{-3}$ & $-2.72994\!\times\!10^{-3}$ & A \\
$\mathbf{2.0721}$ & \multicolumn{2}{c}{\textbf{equal}} & \textbf{indifferent} \\
$3$    & $-5.31529\!\times\!10^{-6}$ & $-4.11151\!\times\!10^{-6}$ & B \\
\bottomrule
\end{tabular}
\end{center}
$$\boxed{\E[u_A] > \E[u_B] \iff \sigma < 2.072} \qquad \checkmark$$

\textbf{Interpretation --- the trade-off is level against dispersion.}

\begin{itemize}[leftmargin=1.3em,itemsep=3pt,topsep=2pt]
\item At $\sigma = 0$ utility is linear, expected utility is just mean consumption, and A
wins by construction ($504 > 402$). \textbf{No inequality aversion at all.}
\item As $\sigma$ rises, the concavity of $u$ increasingly penalises A's wider spread. A's
advantage in the mean is eroded.
\item At $\sigma^{\ast} \approx 2.072$ the two exactly offset: A's $25.4\%$ higher average
consumption is precisely cancelled by its doubled dispersion.
\item Beyond that, \textbf{B wins} --- society would rather be poorer on average than face
A's spread. As $\sigma \to \infty$ only the worst-off individual matters, and since
$c_B(1) = 204 > 108 = c_A(1)$, B wins decisively. That limit is the Rawlsian
difference principle, which Problem 2.10 develops.
\end{itemize}

\textbf{The methodological point.} Whether A or B is the better society is \emph{not} an
empirical question that the data can settle. Both countries' consumption distributions are
fully observed. The ranking depends entirely on $\sigma$ --- a value judgement about how much
weight to give the poor, smuggled in under the guise of a preference parameter. That is both
the strength of the Jones--Klenow approach (the judgement is explicit and you can vary it) and
its limitation (it cannot tell you which $\sigma$ is right).

\end{enumerate}

\begin{tcolorbox}[theorybox,title={Verification code (executed)}]
\footnotesize
\begin{verbatim}
import numpy as np
from scipy.optimize import brentq
j  = np.arange(1, 101)
cA, cB = 100 + 8*j, 200 + 4*j
print(cA.mean(), cB.mean())            # 504.0 402.0

def EU(c, s):
    return np.mean(np.log(c)) if abs(s-1) < 1e-12 else np.mean(c**(1-s)/(1-s))

print(brentq(lambda s: EU(cA, s) - EU(cB, s), 0.05, 3.0))   # 2.072082
\end{verbatim}
\end{tcolorbox}
\end{tcolorbox}

%------- 2.10 -------
% Topic: Risk Aversion, Inequality and the Veil of Ignorance
% Keywords: indifference curves, veil of ignorance, Rawls difference principle, Leontief limit, asymptote
\begin{tcolorbox}[evenbox,title={Problem 2.10 \quad Risk Aversion and the Difference Principle}]
\begin{tcolorbox}[questionbox]
Suppose an individual is trying to evaluate a society from behind the veil of ignorance. He
knows that he can either be rich or poor, with equal probability. Expected utility is
$\E[u(c)] = \dfrac{u(c_{Rich}) + u(c_{Poor})}{2}$ where $u(c) = \dfrac{c^{1-\sigma}}{1-\sigma}$.
Let's now define an indifference curve: for any utility level $\bar U$, the indifference curve
plots all the combinations of $c_{Rich}$ and $c_{Poor}$ such that expected utility is
$\bar U$.\\[3pt]
(a) Take the equation
$\bar U = \dfrac{1}{2}\dfrac{c_{Rich}^{1-\sigma}}{1-\sigma} + \dfrac{1}{2}\dfrac{c_{Poor}^{1-\sigma}}{1-\sigma}$
and solve for $c_{Rich}$.\\
(b) Plot a map of indifference curves for $\sigma = 0.5$ ($\bar U = 2, \sqrt{10}, 4$);
$\sigma = 2$ ($\bar U = -1, -0.4, -0.25$); $\sigma = 5$; $\sigma = 11$. In all cases put
$c_{Poor}$ on the horizontal axis and $c_{Rich}$ on the vertical axis, both ranging over
$[0,10]$. Note that whenever you set $c_{Poor}$ to the left of
$\left[2\bar U(1-\sigma)\right]^{\frac{1}{1-\sigma}}$ you should set $c_{Rich} = \infty$.\\
(c) Notice that as $\sigma$ becomes large, the indifference curves start to look like right
angles. Explain how this relates to Rawls's difference principle.
\end{tcolorbox}

\textbf{Solution:}

\begin{enumerate}

\item \textbf{Solve for $c_{Rich}$.} Starting from
$$\bar U = \frac{1}{2}\cdot\frac{c_{Rich}^{1-\sigma}}{1-\sigma}
+ \frac{1}{2}\cdot\frac{c_{Poor}^{1-\sigma}}{1-\sigma}$$
multiply by $2(1-\sigma)$:
$$2\bar U(1-\sigma) = c_{Rich}^{1-\sigma} + c_{Poor}^{1-\sigma}$$
$$\boxed{\;c_{Rich} = \left[\,2\bar U(1-\sigma) - c_{Poor}^{\,1-\sigma}\,\right]^{\frac{1}{1-\sigma}}\;}$$

\textbf{Where the asymptote comes from.} The bracket must stay positive. For $\sigma < 1$
we need $c_{Poor}^{1-\sigma} < 2\bar U(1-\sigma)$; for $\sigma > 1$ both $1-\sigma$ and the
utility values are negative and the inequality flips. Either way there is a critical value
$$c_{Poor}^{\min} = \left[2\bar U(1-\sigma)\right]^{\frac{1}{1-\sigma}}$$
at which $c_{Rich} \to \infty$: this is precisely the caveat the problem states. \textbf{Below
that level of $c_{Poor}$, no amount of consumption given to the rich can reach utility
$\bar U$} --- the poor person's misery cannot be compensated.

\item \textbf{The plots.} The code below reproduces all four panels. What to look for:

\begin{itemize}[leftmargin=1.3em,itemsep=3pt,topsep=2pt]
\item \textbf{$\sigma = 0.5$:} gently bowed curves, close to straight lines. The vertical
asymptote sits at $c_{Poor}^{\min} = (\bar U)^2$, so at $\bar U = 2$ it is at
$c_{Poor} = 4$. Substitution between rich and poor is easy: taking a little from the
poor is readily offset by giving to the rich.
\item \textbf{$\sigma = 2$:} noticeably more curved; the asymptote is at
$c_{Poor}^{\min} = -1/(2\bar U)$, e.g.\ $0.5$ for $\bar U = -1$.
\item \textbf{$\sigma = 5$ and $\sigma = 11$:} the curves collapse towards \textbf{right
angles}. Over most of the range the indifference curve is essentially vertical (raising
$c_{Rich}$ barely helps) and then essentially horizontal.
\end{itemize}

\item \textbf{Why right angles mean Rawls.}

Right-angled (L-shaped) indifference curves are \textbf{Leontief} preferences,
$U = \min\{c_{Poor}, c_{Rich}\}$ up to a monotone transformation. For such preferences only
the \emph{smaller} argument matters: at the kink, increasing the larger one leaves utility
unchanged.

Formally, the marginal rate of substitution along the curve is
$$\MRS = -\dd{c_{Rich}}{c_{Poor}}
= \left(\frac{c_{Rich}}{c_{Poor}}\right)^{\sigma}$$
Take any point with $c_{Rich} > c_{Poor}$, so the ratio exceeds 1. Then
$$\sigma \to \infty \;\Longrightarrow\; \left(\frac{c_{Rich}}{c_{Poor}}\right)^{\sigma} \to \infty$$
The curve becomes vertical everywhere off the $45^\circ$ line. Interpretation: \textbf{the
planner would give up an unbounded amount of the rich person's consumption for one extra unit
for the poor person.}

That is exactly \textbf{Rawls's difference principle}: social arrangements should be judged
by the position of the \emph{worst-off} member, and an inequality is justified only if it
improves the position of the worst off. In this framework Rawls is not a separate ethical
theory bolted on --- it is the \textbf{limiting case of expected-utility maximisation behind
the veil of ignorance as risk aversion goes to infinity}.

\medskip
\textbf{Why this matters for Chapter 2.} It locates Jones and Klenow's choice of
$\sigma$ on a spectrum of ethical positions rather than treating it as a technical
calibration:
\begin{center}
\begin{tabular}{ll}
\toprule
$\sigma = 0$ & pure utilitarian --- maximise \emph{mean} consumption, inequality irrelevant \\
$\sigma = 1$ & log utility --- Jones \& Klenow's baseline; maximise the geometric mean \\
$\sigma \to \infty$ & Rawlsian --- maximise the consumption of the worst-off \\
\bottomrule
\end{tabular}
\end{center}
Choosing $\sigma$ \emph{is} choosing a social welfare criterion. Problem 2.8 showed this has
real bite: it flipped the ranking of two countries at $\sigma^{\ast} \approx 2.07$.

\end{enumerate}

\begin{tcolorbox}[theorybox,title={Plotting code}]
\footnotesize
\begin{verbatim}
import numpy as np, matplotlib.pyplot as plt

def c_rich(c_poor, Ubar, sigma):
    br = 2*Ubar*(1-sigma) - c_poor**(1-sigma)
    with np.errstate(invalid="ignore"):
        out = np.where(br > 0, np.abs(br)**(1/(1-sigma)), np.inf)
    return out

panels = [(0.5, [2, np.sqrt(10), 4]), (2, [-1, -0.4, -0.25]),
          (5, [-0.25, -4/54, -4.0**-5]), (11, [-0.1, -4.0**10/10**11, -4.0**-10/10])]

fig, axes = plt.subplots(2, 2, figsize=(10, 9))
cp = np.linspace(0.01, 10, 2000)
for ax, (sig, Us) in zip(axes.ravel(), panels):
    for U in Us:
        ax.plot(cp, c_rich(cp, U, sig), lw=1.4, label=f"Ubar={U:.3g}")
    ax.set_xlim(0, 10); ax.set_ylim(0, 10)
    ax.plot([0, 10], [0, 10], ls=":", c="gray", lw=.8)     # 45-degree line
    ax.set_title(f"sigma = {sig}"); ax.set_xlabel("c_Poor")
    ax.set_ylabel("c_Rich"); ax.legend(fontsize=7)
plt.tight_layout()
\end{verbatim}
\end{tcolorbox}
\end{tcolorbox}

%======================================================================
\section{Chapter 3 \quad Basic Facts about Economic Growth}
%======================================================================

\begin{tcolorbox}[theorybox,title={Chapter 3 Overview}]
\textbf{What this chapter is for.} Before building a model of growth, establish the facts it
has to explain. All three exercises are \textbf{empirical} --- they ask you to go and look.

\medskip
\textbf{The very long run.} Living standards were roughly flat for millennia and then took
off around 1800. Compounding does the work: $1.5\%$ a year sounds trivial and multiplies
income by $4.4$ over a century.

\medskip
\textbf{The Kaldor facts.} Over long horizons in developed economies: (1) $Y/L$ grows at a
roughly constant rate; (2) $K/L$ grows at a roughly constant rate; (3) $K/Y$ is roughly
constant; (4) the return on capital is roughly constant; (5) factor shares are roughly
constant. Facts 2 and 3 together imply fact 4, but it is stated separately because the
behaviour of $r$ matters independently. These five facts are the target the Solow model of
Chapter 4 is built to hit.

\medskip
\textbf{Across countries.} Initially-rich countries grow at similar, middling rates.
Initially-poor countries show \emph{enormous} dispersion --- Korea and Botswana race ahead
while Congo and Madagascar fall further behind. So there is \textbf{no absolute convergence}
in the world sample. Whether there is convergence \emph{within} more homogeneous groups is
exactly what Exercise 3.3 asks.

\medskip
\textbf{Reading growth data.} Always plot income in \textbf{logs}: then a constant growth
rate is a straight line and the \emph{slope} is the growth rate. Comparing slopes across
countries is comparing growth; comparing heights is comparing levels.
\end{tcolorbox}

\begin{tcolorbox}[databox,title={All three exercises in this chapter are data exercises}]
Exercises 3.1--3.3 require downloading the Maddison Project Database and the Penn World
Table. No numerical answers are asserted below. For each one I give the exact procedure, the
pattern you should expect to find and why, and runnable code. Where I give an approximate
historical figure it is flagged as approximate --- the precise answer depends on the data
release and the PPP benchmark year, which is itself part of the lesson.
\end{tcolorbox}

%---- Topic: Long-Run Growth and Historical Comparisons ----
\subsection{Long-Run Growth and Historical Comparisons}

%------- 3.1 -------
% Topic: Long-Run Growth and Historical Comparisons
% Keywords: Maddison Project, historical GDP per capita, compounding, PPP benchmark, level comparisons
\begin{tcolorbox}[oddbox,title={Problem 3.1 \quad The Past is a Foreign Country}]
\begin{tcolorbox}[questionbox]
Find the data compiled by Bolt et al.\ (2018) at
\url{https://www.rug.nl/ggdc/historicaldevelopment/maddison/releases/maddison-project-database-2018}.
Starting from 1800, look up the US GDP per capita at intervals of one decade. For each of
these points in time, find the country that currently has the closest GDP per capita to the
past US level. When did the US become richer than present-day India, present-day China, and
present-day Portugal?
\end{tcolorbox}

\textbf{Solution:}

\begin{enumerate}

\item \textbf{Procedure.} Download \texttt{mpd2018.xlsx}, sheet \texttt{Full data}. The
variable you want is \texttt{cgdppc} (real GDP per capita in 2011 US dollars, PPP-adjusted),
which is the one constructed for \emph{cross-country level} comparisons --- as opposed to
\texttt{rgdpnapc}, which is built for comparisons \emph{over time within} a country. Getting
this choice right is the whole methodological point of the exercise: you are comparing the
past US to the present of other countries, which is a level comparison across both time and
space.

Then, for each decade $t \in \{1800, 1810, \ldots\}$, find the country $j$ minimising
$\left|y^{US}_t - y^{j}_{\text{latest}}\right|$.

\item \textbf{Approximate answers --- verify these against the data yourself.}

The threshold years are found by locating where the US series crosses each country's current
level:

\begin{center}
\begin{tabular}{lll}
\toprule
\textbf{Present-day country} & \textbf{Approx.\ current level} & \textbf{US passed it around} \\
\midrule
India    & \$6,000--7,000   & the \textbf{1870s--1880s} \\
China    & \$13,000--16,000 & the \textbf{1900s--1910s} \\
Portugal & \$27,000--30,000 & the \textbf{1940s--1950s} \\
\bottomrule
\end{tabular}
\end{center}

\begin{tcolorbox}[databox,title={Treat these as orders of magnitude, not answers}]
China's level in particular moves fast --- its GDP per capita roughly doubled in the decade
before the 2018 release --- so the crossing year is sensitive to the vintage. Report the year
\emph{and} the release you used.
\end{tcolorbox}

\item \textbf{What the exercise is teaching.} Three things, in increasing order of
importance:

\begin{itemize}[leftmargin=1.3em,itemsep=3pt,topsep=2pt]
\item \textbf{The magnitude of modern growth.} The US of the 1880s --- railroads,
      telegraphs, mass steel --- had roughly the material living standard of India today.
      Two centuries of compounding at rates that look small in any given year produced a
      gap of more than an order of magnitude.
\item \textbf{Time and space are interchangeable measures of development.} ``How rich is
      India?'' and ``how far back is the US equivalent?'' are the same question. This is
      what makes the title apt.
\item \textbf{The comparison is cruder than it looks.} An 1880 American and a 2018 Indian
      with the same measured real income do not consume remotely the same bundle. The Indian
      has antibiotics, vaccines, mobile telephony and electricity; the American had none of
      them. Chained price indices (Chapter 1) cannot handle goods that did not exist ---
      the \textbf{new-goods bias} --- so this kind of comparison almost certainly
      \emph{understates} how much better off the present-day poor country is. That caveat is
      worth stating in your answer.
\end{itemize}

\end{enumerate}

\begin{tcolorbox}[theorybox,title={Code}]
\footnotesize
\begin{verbatim}
import pandas as pd

mpd = pd.read_excel("mpd2018.xlsx", sheet_name="Full data")
us  = mpd[mpd.countrycode == "USA"].set_index("year")["cgdppc"]

latest = (mpd.dropna(subset=["cgdppc"]).sort_values("year")
             .groupby("country").tail(1).set_index("country")["cgdppc"])

for decade in range(1800, 2011, 10):
    if decade in us.index and pd.notna(us[decade]):
        match = (latest - us[decade]).abs().idxmin()
        print(f"{decade}  US = {us[decade]:8,.0f}   closest today: {match} "
              f"({latest[match]:,.0f})")

for target in ["India", "China", "Portugal"]:
    lvl = latest[target]
    crossed = us[us > lvl]
    print(f"US passed present-day {target} (${lvl:,.0f}) in {crossed.index.min()}")
\end{verbatim}
\end{tcolorbox}
\end{tcolorbox}

%---- Topic: The Kaldor Facts ----
\subsection{The Kaldor Facts}

%------- 3.2 -------
% Topic: The Kaldor Facts
% Keywords: Kaldor facts, Penn World Table, capital-output ratio, labor share, balanced growth
\begin{tcolorbox}[evenbox,title={Problem 3.2 \quad The Kaldor Facts in Other Countries}]
\begin{tcolorbox}[questionbox]
Look up GDP accounts for some countries other than the US (a good source is
\url{https://www.rug.nl/ggdc/productivity/pwt/}). Reproduce Figures 3.2.1--3.2.3 using data
from those countries (try a rich country, a middle-income country and a poor country). Do the
Kaldor facts seem to hold for those countries as well?
\end{tcolorbox}

\textbf{Solution:}

\begin{enumerate}

\item \textbf{What to build, and from which PWT variables.}

\begin{center}
\begin{tabular}{lll}
\toprule
\textbf{Kaldor fact} & \textbf{Plot} & \textbf{PWT 10.x variables} \\
\midrule
1. $Y/L$ grows steadily & $\ln(Y/L)$ vs year & \texttt{rgdpna} / \texttt{emp} \\
2. $K/L$ grows steadily & $\ln(K/L)$ vs year & \texttt{rnna} / \texttt{emp} \\
3. $K/Y$ constant       & $K/Y$ vs year      & \texttt{rnna} / \texttt{rgdpna} \\
5. Factor shares constant & labour share vs year & \texttt{labsh} \\
\bottomrule
\end{tabular}
\end{center}

Use \texttt{rgdpna} (national-accounts real GDP) for \emph{within-country over time}, not
\texttt{cgdpo}. Fact 4 (the return on capital) is not directly observed; it can be inferred
from \texttt{labsh} and $K/Y$ via $r + \delta = \alpha \frac{Y}{K}$.

\item \textbf{Suggested countries and what to expect.}

\begin{itemize}[leftmargin=1.3em,itemsep=4pt,topsep=2pt]
\item \textbf{Rich (e.g.\ France, Japan, Canada).} The facts should hold reasonably well
      --- this is the sample Kaldor generalised from. Expect a caveat: Japan's $K/Y$ rose
      substantially through its high-investment era and its growth rate is clearly \emph{not}
      constant (rapid to 1990, then flat), so ``constant growth'' holds only piecewise.
      Labour shares in most rich countries have \textbf{drifted down} since around 1980 ---
      a well-documented violation of fact 5 and an active research area.
\item \textbf{Middle-income (e.g.\ Brazil, Mexico, Thailand).} Expect the facts to hold
      \emph{much} worse. Brazil is a good case precisely because it fails so visibly: rapid
      growth to about 1980, then the lost decade, then stagnation. $\ln(Y/L)$ is not a
      straight line by any reading. $K/Y$ moves substantially.
\item \textbf{Poor (e.g.\ Madagascar, Congo, Malawi).} Expect the facts to fail badly.
      $\ln(Y/L)$ may be \emph{flat or falling} over decades. Data quality is also weakest
      here --- a point worth making explicitly rather than treating the series as gospel.
\end{itemize}

\item \textbf{The answer to the question as asked: the Kaldor facts are properties of
economies on a balanced growth path, not laws of nature.}

They hold approximately for rich countries over long horizons --- which is exactly the domain
Kaldor proposed them for --- and they \textbf{fail systematically for countries that are not
on such a path.} That is not a defect of the data; it is informative. A country in the middle
of a growth miracle or a growth disaster is \emph{in transition}, and the Solow model
(Chapter 4) predicts that transitional dynamics look precisely like this: $K/Y$ rising as a
country accumulates towards its steady state, $Y/L$ growth decelerating as diminishing
returns bite.

\textbf{So the honest conclusion has two parts:} the facts describe steady states well, and
most of the world is not in a steady state. Chapter 5 takes this up directly when it asks
whether capital accumulation alone can account for cross-country income differences.

\end{enumerate}

\begin{tcolorbox}[theorybox,title={Code}]
\footnotesize
\begin{verbatim}
import pandas as pd, numpy as np, matplotlib.pyplot as plt

pwt = pd.read_excel("pwt1001.xlsx", sheet_name="Data")
countries = ["France", "Brazil", "Madagascar"]     # rich / middle / poor

fig, axes = plt.subplots(3, 3, figsize=(13, 10), sharex=True)
for row, ctry in enumerate(countries):
    d = pwt[pwt.country == ctry].set_index("year")
    y_l = d.rgdpna / d.emp                          # fact 1
    k_l = d.rnna   / d.emp                          # fact 2
    k_y = d.rnna   / d.rgdpna                       # fact 3

    axes[row,0].plot(np.log(y_l));  axes[row,0].set_ylabel(f"{ctry}\nlog(Y/L)")
    axes[row,1].plot(np.log(k_l));  axes[row,1].set_ylabel("log(K/L)")
    axes[row,2].plot(k_y);          axes[row,2].set_ylabel("K/Y")
    print(ctry, "labour share: first", d.labsh.dropna().iloc[0].round(3),
                "last", d.labsh.dropna().iloc[-1].round(3))
plt.tight_layout()
\end{verbatim}
\end{tcolorbox}
\end{tcolorbox}

%---- Topic: Growth Across Countries and Convergence ----
\subsection{Growth Across Countries and Convergence}

%------- 3.3 -------
% Topic: Growth Across Countries and Convergence
% Keywords: absolute convergence, conditional convergence, club convergence, scatterplot, regional samples
\begin{tcolorbox}[oddbox,title={Problem 3.3 \quad Within-Region Convergence}]
\begin{tcolorbox}[questionbox]
Look up cross-country GDP data (a good source is
\url{https://www.rug.nl/ggdc/productivity/pwt/}). For each country, compute average
GDP-per-capita growth 1960--2014. Produce a scatterplot of growth against initial GDP per
capita like Figure 3.3.1 but separately for countries in each of three regions of the world:
Europe, Latin America and Africa. Is it the case in any of these regions that initially-poor
countries have grown faster?
\end{tcolorbox}

\textbf{Solution:}

\begin{tcolorbox}[theorybox,title={This is question 3 of Lista de Exerc\'icios 1 (2 pts)}]
The Lista assigns this exercise directly. Note that it is a \textbf{data exercise} --- you
will need to download the Penn World Table and actually produce three scatterplots. Budget
time for it; it is not a paper-and-pencil question. Deliver the three plots plus the fitted
slopes.
\end{tcolorbox}

\begin{enumerate}

\item \textbf{Procedure.} For each country with data in both years, compute the average
annual growth rate
$$g_i = \left(\frac{y_{i,2014}}{y_{i,1960}}\right)^{1/54} - 1$$
and plot $g_i$ against $\ln y_{i,1960}$ (logs on the horizontal axis --- this is the standard
convention and makes the convergence slope interpretable). Fit a line to each regional
subsample. \textbf{A negative slope is evidence of convergence.}

\item \textbf{What you should find, and why.}

\begin{center}
\begin{tabular}{lll}
\toprule
\textbf{Region} & \textbf{Expected slope} & \textbf{Reading} \\
\midrule
\textbf{Europe} & clearly \textbf{negative} & strong convergence \\
\textbf{Latin America} & mildly negative or flat & weak / no convergence \\
\textbf{Africa} & flat or \textbf{positive} & no convergence; possible divergence \\
\bottomrule
\end{tabular}
\end{center}

\begin{itemize}[leftmargin=1.3em,itemsep=4pt,topsep=2pt]
\item \textbf{Europe converges, and convincingly.} In 1960 Portugal, Spain, Greece and
      Ireland were far poorer than Germany, France and the UK; by 2014 the gap had narrowed
      dramatically. Ireland is the extreme case. Add the post-1990 Eastern European countries
      and the negative slope is stronger still.
\item \textbf{Latin America shows little.} Countries started at broadly similar middle-income
      levels and mostly stayed there --- the region as a whole failed to converge on the US.
      Within-region dispersion did not narrow much; Venezuela moved sharply in the wrong
      direction.
\item \textbf{Africa shows none, and possibly the reverse.} Botswana and Mauritius grew
      rapidly from low bases while Congo, Madagascar, Niger and the Central African Republic
      stagnated or shrank. Initial income predicts almost nothing.
\end{itemize}

\item \textbf{Interpretation --- this is the distinction between absolute and conditional
convergence.} The exercise is engineered to make the point that \emph{the world sample shows
no absolute convergence, but sub-samples can}.

The Solow model predicts that each country converges to \textbf{its own} steady state $k^\ast$,
which depends on its $s$, $n$, $\delta$ and technology. That gives \textbf{conditional
convergence}: countries approach \emph{their own} target, faster the further away they are.
It gives \emph{absolute} convergence only if all countries share the same steady state.

\begin{itemize}[leftmargin=1.3em,itemsep=3pt,topsep=2pt]
\item Within \textbf{Europe} --- similar institutions, similar saving and demographic rates,
      free movement of capital, technology and labour, and for many the disciplining effect
      of EU accession --- the assumption of a \emph{common} steady state is roughly
      defensible. So conditional convergence shows up as absolute convergence.
\item Within \textbf{Africa}, steady states differ enormously: institutions, conflict, health,
      schooling, geography and policy vary far more than initial income does. Countries are
      converging, but to widely different targets, so nothing shows up in a plot against
      initial income alone.
\end{itemize}

This is the phenomenon of \textbf{convergence clubs}. Homogeneous groups converge internally;
the world as a whole does not.

\item \textbf{The caveat that earns marks.} Restricting attention to a region conditions on
\emph{observable} regional membership, but it is still not the same as controlling for $s$
and $n$ directly. A proper conditional-convergence test regresses
$$g_i = \beta_0 + \beta_1 \ln y_{i,1960} + \beta_2 \ln s_i + \beta_3 \ln(n_i + g + \delta)
+ \varepsilon_i$$
and asks whether $\beta_1 < 0$. The regional split is a crude, transparent proxy for that.
Also beware \textbf{selection}: requiring data in both 1960 and 2014 drops exactly the
countries whose statistical systems collapsed --- typically the worst performers --- which
biases African growth \emph{upward} and understates divergence.

\end{enumerate}

\begin{tcolorbox}[theorybox,title={Code}]
\footnotesize
\begin{verbatim}
import pandas as pd, numpy as np, matplotlib.pyplot as plt

pwt = pd.read_excel("pwt1001.xlsx", sheet_name="Data")
pwt["y"] = pwt.rgdpna / pwt.pop

w = (pwt[pwt.year.isin([1960, 2014])]
       .pivot(index="country", columns="year", values="y").dropna())
w["g"]    = (w[2014] / w[1960]) ** (1/54) - 1
w["ly60"] = np.log(w[1960])

regions = {
 "Europe":  ["France","Germany","Italy","Spain","Portugal","Greece","Ireland",
             "Netherlands","Belgium","Austria","Denmark","Sweden","Norway",
             "Finland","Switzerland","United Kingdom"],
 "Latin America": ["Brazil","Argentina","Chile","Colombia","Peru","Mexico",
             "Uruguay","Venezuela (Bolivarian Republic of)","Ecuador","Bolivia (Plurinational State of)",
             "Paraguay","Costa Rica","Panama","Guatemala","Honduras"],
 "Africa":  ["Botswana","Kenya","Nigeria","Ghana","Senegal","Cameroon","Niger",
             "Malawi","Zambia","Zimbabwe","Mozambique","Madagascar","Mauritius",
             "South Africa","Central African Republic"],
}

fig, axes = plt.subplots(1, 3, figsize=(15, 4.5), sharey=True)
for ax, (name, members) in zip(axes, regions.items()):
    d = w[w.index.isin(members)]
    ax.scatter(d.ly60, d.g)
    b, a = np.polyfit(d.ly60, d.g, 1)                 # slope, intercept
    xs = np.linspace(d.ly60.min(), d.ly60.max(), 10)
    ax.plot(xs, a + b*xs, "r-")
    ax.set_title(f"{name}: slope = {b:+.4f}")
    ax.set_xlabel("log GDP per capita, 1960")
axes[0].set_ylabel("average growth 1960-2014")
plt.tight_layout()
\end{verbatim}
\end{tcolorbox}
\end{tcolorbox}

%======================================================================
\section{Chapter 4 \quad The Solow Growth Model}
%======================================================================

\begin{tcolorbox}[theorybox,title={Chapter 4 Overview}]
\textbf{The assumptions on $F$} (numbered as in \S4.1, and you will need them by number in
Problem 4.4):
\begin{itemize}[leftmargin=1.3em,itemsep=2pt,topsep=2pt]
\item \textbf{4.1 Constant Returns to Scale:} $F(\lambda K, \lambda L) = \lambda F(K,L)$.
\item \textbf{4.2 Positive Marginal Product:} $F_K > 0$, $F_L > 0$.
\item \textbf{4.3 Diminishing Marginal Product:} $F_{KK} < 0$, $F_{LL} < 0$.
\item \textbf{4.4 Inada Conditions:} $\lim_{K\to 0}F_K = \infty$, $\lim_{K\to\infty}F_K = 0$.
\item \textbf{4.5 Population:} $L_{t+1} = (1+n)L_t$, with $n$ exogenous.
\end{itemize}
CRS is what lets us work per worker: $y = f(k)$ with $k = K/L$.

\medskip
\textbf{Mechanics (\S4.2).} $\;k_{t+1} = \dfrac{(1-\delta)k_t + s f(k_t)}{1+n}$, with steady
state $s f(k\stst) = (n+\delta)k\stst$. It is globally stable, and Inada guarantees it exists
and is unique. \textbf{The steady state depends on $s$, $n$, $\delta$ --- never on the initial
$k_0$.} This single fact answers Problems 4.1 and 4.2.

\medskip
\textbf{Golden Rule (\S4.3).} $f'(k_{gold}) = n+\delta$; with Cobb--Douglas, $s_{gold} = \alpha$.

\medskip
\textbf{Factor markets (\S4.4).} The results Problem 4.3 runs on:
$$r^K = F_K(k,1), \qquad w = F_L(1, 1/k), \qquad \boxed{r = r^K - \delta} \;\;\text{(eq.\ 4.4.12)}$$
With Cobb--Douglas $Y = K^\alpha L^{1-\alpha}$:
$$r^K = \alpha k^{\alpha-1}, \qquad w = (1-\alpha)k^{\alpha}$$
so \textbf{high $k$ means high wages and a low return on capital}, and vice versa. Profits are
zero (Proposition 4.2) --- all output goes to workers or to capital owners.

\medskip
\textbf{Technological progress (\S4.5).} Only with $A_{t+1}=(1+g)A_t$ does the model deliver
long-run growth in $y$, and it does so \emph{by assumption}. Without it, the Solow model
predicts no long-run growth in GDP per capita at all.
\end{tcolorbox}

%---- Topic: Comparative Statics --- Shocks to Capital and Labour ----
\subsection{Comparative Statics --- Shocks to Capital and Labour}

%------- 4.1 -------
% Topic: Comparative Statics --- Shocks to Capital and Labour
% Keywords: steady state independence of initial conditions, transitional dynamics, convergence, factor prices
\begin{tcolorbox}[oddbox,title={Problem 4.1 \quad An Earthquake}]
\begin{tcolorbox}[questionbox]
Suppose an economy behaves according to the Solow growth model. It starts out at $t=0$ at a
steady state, with no technological progress and no population growth. Suppose an earthquake
destroys half the capital stock. As a consequence of this, what would happen in the short run
(i.e.\ immediately) and in the long run (i.e.\ once the economy reaches steady state) to:
(a) GDP, (b) GDP per capita, (c) wages, (d) interest rates.
\end{tcolorbox}

\textbf{Solution:}

\begin{tcolorbox}[theorybox,title={This is question 4 of Lista de Exerc\'icios 1 (3 pts)}]
The Lista uses this exercise almost verbatim, then adds a part (b) in which the population
growth rate is \emph{also} halved because people emigrate. That extension changes the
long-run answer --- see the note at the end of this solution.
\end{tcolorbox}

\emph{Setup.} $n = 0$, $g = 0$. The initial steady state satisfies $s f(k\stst) = \delta k\stst$.
The earthquake sets $K \to K/2$ with $L$ unchanged, so
$$k_0 = \frac{k\stst}{2}$$

\textbf{The organising principle: $k\stst$ depends only on $s$, $n$, $\delta$ and the
production function --- not on the initial condition.} The earthquake changes the initial
condition and nothing else. Therefore \textbf{every long-run answer is ``returns to exactly
where it was''}, and all the action is in the transition.

\begin{center}
\begin{tabular}{llll}
\toprule
\textbf{Variable} & \textbf{Short run} & \textbf{Long run} & \textbf{Because} \\
\midrule
(a) GDP          & $\downarrow$ & back to original & $Y = F(K,L)$, $K$ halved \\
(b) GDP per capita & $\downarrow$ & back to original & $y = f(k)$, $k$ halved \\
(c) Wages        & $\downarrow$ & back to original & $w = (1-\alpha)k^\alpha$ \\
(d) Interest rate & $\uparrow$   & back to original & $r = \alpha k^{\alpha-1} - \delta$ \\
\bottomrule
\end{tabular}
\end{center}

\begin{enumerate}

\item \textbf{GDP falls immediately, but by \emph{less} than half.} With Cobb--Douglas,
$$\frac{Y_{new}}{Y_{old}} = \left(\frac{K/2}{K}\right)^{\alpha} = \left(\tfrac12\right)^{\alpha}$$
With $\alpha = 1/3$ this is $0.5^{1/3} = 0.794$, so GDP falls about $21\%$ --- not $50\%$.
The reason is \textbf{diminishing returns} (Assumption 4.3): the capital destroyed was the
\emph{least} productive capital at the margin, and labour is unaffected.

In the long run $k$ returns to $k\stst$ and $Y$ returns to its original level.

\item \textbf{GDP per capita: same as (a).} Since $L$ is unchanged, $y = Y/L$ moves exactly
in proportion to $Y$. Falls by $1 - 0.5^\alpha$ on impact, returns fully in the long run.

\item \textbf{Wages fall.} $w = F_L = (1-\alpha)k^{\alpha}$ is \emph{increasing} in $k$.
Workers have less capital to work with, so labour's marginal product falls:
$$\frac{w_{new}}{w_{old}} = \left(\tfrac12\right)^{\alpha}$$
the same proportional drop as output --- because with Cobb--Douglas the labour share is
constant at $1-\alpha$, so $w$ and $y$ always move together. Returns to the original level in
the long run.

\item \textbf{Interest rates rise.} $r = r^K - \delta = \alpha k^{\alpha-1} - \delta$ is
\emph{decreasing} in $k$ (Assumption 4.3 again). Capital is now scarce, so its marginal
product is high:
$$r^K_{new} = \alpha\left(\tfrac{k\stst}{2}\right)^{\alpha-1}
= 2^{\,1-\alpha}\,\alpha (k\stst)^{\alpha-1} = 2^{\,1-\alpha}\, r^K_{old}$$
With $\alpha = 1/3$ the rental rate rises by a factor $2^{2/3} = 1.587$, i.e.\ about $59\%$.
Returns to the original level in the long run.

\end{enumerate}

\textbf{The transition, which is the economically interesting part.} At $k_0 = k\stst/2$ we
have $k < k\stst$, so
$$s f(k_0) > \delta k_0$$
--- investment exceeds depreciation, and $k$ \emph{rises} back towards $k\stst$. During this
period the economy grows \textbf{faster than its long-run rate} (which is zero here). Growth
is fastest immediately after the shock and slows as $k$ approaches $k\stst$.

\textbf{Why this matters empirically.} This is the Solow model's explanation of the post-war
German and Japanese ``growth miracles''. Wartime destruction pushed $k$ far below $k\stst$;
the rapid growth that followed was not a change in those countries' fundamentals but
\textbf{convergence back to an unchanged steady state}. It is also the model's prediction
that natural disasters have no permanent effect on income --- a prediction that holds up
reasonably well for rich countries with intact institutions, and much less well where the
disaster damages institutions or human capital, neither of which is in the model.

\begin{tcolorbox}[theorybox,title={Extension in Lista 1: what if $n$ halves too?}]
The Lista adds that people emigrate, halving the population growth rate $n$. Now the shock
hits a \emph{parameter}, not just the initial condition, so the long-run answer changes:
$$k\stst = \left(\frac{s}{n+\delta}\right)^{\frac{1}{1-\alpha}}
\quad\Longrightarrow\quad n \downarrow \;\Rightarrow\; k\stst \uparrow$$
\textbf{GDP per capita ends up \emph{higher} than before the earthquake}, along with wages;
the interest rate ends up \emph{lower}. Total GDP is ambiguous in level terms and grows more
slowly ($n$ is lower). Note the exercise says $n=0$ initially --- if you read ``halved'' as
applying to a strictly positive $n$, say so explicitly; if $n = 0$ then halving it changes
nothing and the answer reverts to the base case. \textbf{State which reading you adopt} ---
this ambiguity is the kind of thing the rubric in \texttt{rules/00\_estilo\_avaliacao.md}
treats as a residual risk.
\end{tcolorbox}
\end{tcolorbox}

%------- 4.2 -------
% Topic: Comparative Statics --- Shocks to Capital and Labour
% Keywords: Black Death, labour scarcity, capital deepening, factor prices, historical evidence
\begin{tcolorbox}[evenbox,title={Problem 4.2 \quad The Black Death}]
\begin{tcolorbox}[questionbox]
In the middle of the 14th century, a plague killed about a third of the population of Europe.
Assume that the economy of Europe was well described by the Solow model. How would the
following variables change in response to the Black Death in the short run (i.e.\
immediately)? (a) GDP, (b) GDP per capita, (c) wages, (d) interest rates.
\end{tcolorbox}

\textbf{Solution:}

\emph{Setup.} $L \to \tfrac{2}{3}L$ with $K$ unchanged. Therefore
$$k = \frac{K}{L} \;\longrightarrow\; \frac{K}{(2/3)L} = \tfrac{3}{2}k$$
\textbf{Capital per worker rises by 50\%.} This problem is the mirror image of 4.1: there
capital was destroyed and $k$ fell; here labour is destroyed and $k$ rises. Every sign flips.

\begin{center}
\begin{tabular}{lll}
\toprule
\textbf{Variable} & \textbf{Short run} & \textbf{With $\alpha = 1/3$} \\
\midrule
(a) GDP            & $\downarrow$ & $\times (2/3)^{1-\alpha} = 0.763$ \quad ($-24\%$) \\
(b) GDP per capita & $\uparrow$   & $\times (3/2)^{\alpha} = 1.145$ \quad ($+15\%$) \\
(c) Wages          & $\uparrow$   & $\times (3/2)^{\alpha} = 1.145$ \quad ($+15\%$) \\
(d) Interest rate  & $\downarrow$ & $r^K \times (3/2)^{\alpha-1} = 0.874$ \quad ($-13\%$) \\
\bottomrule
\end{tabular}
\end{center}

\begin{enumerate}

\item \textbf{GDP falls.} Fewer workers, same capital:
$$\frac{Y_{new}}{Y_{old}} = \left(\frac{2}{3}\right)^{1-\alpha}$$
With $\alpha = 1/3$: $(2/3)^{2/3} = 0.763$, so total output falls about $24\%$. Note it falls
by \emph{less} than the third of the population that died --- diminishing returns again, this
time to labour.

\item \textbf{GDP per capita \emph{rises}.}
$$\frac{y_{new}}{y_{old}} = \left(\frac{k_{new}}{k_{old}}\right)^{\alpha}
= \left(\tfrac{3}{2}\right)^{1/3} = 1.145$$
about $+15\%$. This is the striking result and the point of the exercise: \textbf{the
survivors were, on average, materially better off.} Output fell by less than population did,
so output per head rose.

\item \textbf{Wages rise, by the same $15\%$.} $w = (1-\alpha)k^\alpha$ and $k$ rose. Labour
became the scarce factor and its marginal product rose accordingly.

\textbf{This is the model's most successfully tested prediction.} Real wages in England
roughly doubled over the century following 1348, and the historical record shows exactly the
pattern the model implies: acute labour scarcity, landlords competing for tenants and
labourers, and legislative attempts to suppress the adjustment --- the English Statute of
Labourers of 1351 tried to fix wages at pre-plague levels and largely failed. That legislation
is itself evidence that wages were rising sharply.

\item \textbf{Interest rates fall.} $r^K = \alpha k^{\alpha - 1}$ is decreasing in $k$:
$$\frac{r^K_{new}}{r^K_{old}} = \left(\tfrac{3}{2}\right)^{\alpha-1}
= \left(\tfrac{3}{2}\right)^{-2/3} = 0.874$$
Capital became abundant relative to labour, so its return fell. Since $r = r^K - \delta$, the
interest rate falls by the same absolute amount as the rental rate. Land rents fell for the
same reason --- another well-documented feature of the post-plague period.

\end{enumerate}

\begin{tcolorbox}[theorybox,title={The long run --- and why history did not stay there}]
The question asks only about the short run, but the long run is worth a sentence. With $n$
unchanged, population recovers, $k$ falls back to $k\stst$, and \textbf{wages and income per
head return to their original levels.} The gains to survivors were transitory.

That is roughly what happened historically over the following two centuries --- and it is
precisely the \textbf{Malthusian} mechanism of Problem 4.5, where population growth responds
endogenously to living standards. In the Solow model with exogenous $n$ the return to
$k\stst$ is mechanical; in the Malthusian model it is \emph{driven} by the high living
standards themselves. Problem 4.5 makes that explicit.
\end{tcolorbox}

\begin{tcolorbox}[theorybox,title={Comparing 4.1 and 4.2 --- the pairing is deliberate}]
\begin{center}
\begin{tabular}{lcc}
\toprule
& \textbf{4.1 Earthquake ($K\downarrow$)} & \textbf{4.2 Plague ($L\downarrow$)} \\
\midrule
$k = K/L$          & falls & \textbf{rises} \\
GDP                & falls & falls \\
GDP per capita     & falls & \textbf{rises} \\
Wages              & falls & \textbf{rises} \\
Interest rate      & rises & \textbf{falls} \\
\bottomrule
\end{tabular}
\end{center}
Both shocks reduce total output. They have \textbf{opposite} effects on everything per-capita
and on both factor prices, because they move $k$ in opposite directions. If you can state
\emph{why} --- factor prices are determined by $k$ alone, via $w = (1-\alpha)k^\alpha$ and
$r^K = \alpha k^{\alpha-1}$ --- you have understood \S4.4.
\end{tcolorbox}
\end{tcolorbox}

%---- Topic: Factor Markets and Factor Prices ----
\subsection{Factor Markets and Factor Prices}

%------- 4.3 -------
% Topic: Factor Markets and Factor Prices
% Keywords: marginal products, wages, rental rate, factor mobility, Jensen inequality, gains from integration
\begin{tcolorbox}[oddbox,title={Problem 4.3 \quad Korean Unification}]
\begin{tcolorbox}[questionbox]
Suppose both North and South Korea have the same technology level (!?), which doesn't grow,
and also have constant population. Their respective populations and capital stocks are
$L_{South} = 10$, $K_{South} = 2430$, $L_{North} = 5$, $K_{North} = 160$. The production
function is $Y = K^{\alpha}L^{1-\alpha}$ with $\alpha = 0.4$. The depreciation rate is
$\delta = 0.08$.\\[3pt]
(a) Compute GDP per capita in each country.\\
(b) Compute wages and interest rates in each country.\\
(c) Suppose North and South Korea unify. What is the new $\frac{K}{L}$ in the unified
country? What is GDP per capita in the unified country?\\
(d) Compute wages and interest rates in the unified country. Interpret.\\
(e) Will people and/or machines physically move between the North and the South? In what
direction?
\end{tcolorbox}

\textbf{Solution:}

\emph{Formulas.} With $Y = K^{0.4}L^{0.6}$:
$$y = f(k) = k^{0.4}, \qquad w = F_L = (1-\alpha)k^{\alpha} = 0.6\,k^{0.4},$$
$$r^K = F_K = \alpha k^{\alpha-1} = 0.4\,k^{-0.6}, \qquad r = r^K - \delta = r^K - 0.08$$

\begin{enumerate}

\item \textbf{GDP per capita.} The numbers are chosen to come out exactly:
$$k_S = \frac{2430}{10} = 243 = 3^5 \;\Longrightarrow\; y_S = 243^{0.4} = (3^5)^{0.4} = 3^2
= \boxed{9}$$
$$k_N = \frac{160}{5} = 32 = 2^5 \;\Longrightarrow\; y_N = 32^{0.4} = (2^5)^{0.4} = 2^2
= \boxed{4}$$
\checkmark\ Total output: $Y_S = 9 \times 10 = 90$ and $Y_N = 4 \times 5 = 20$.

The South is $2.25$ times richer per head, purely because it has $7.6$ times more capital per
worker.

\item \textbf{Wages and interest rates.}

\begin{center}
\begin{tabular}{lrrrr}
\toprule
& $k$ & $w = 0.6\,k^{0.4}$ & $r^K = 0.4\,k^{-0.6}$ & $r = r^K - 0.08$ \\
\midrule
South & 243 & $0.6 \times 9 = \mathbf{5.4}$ & $0.4/27 = \mathbf{0.01481}$ & $\mathbf{-0.0652}$ \\
North & 32  & $0.6 \times 4 = \mathbf{2.4}$ & $0.4/8 \;= \mathbf{0.05000}$ & $\mathbf{-0.0300}$ \\
\bottomrule
\end{tabular}
\end{center}
\checkmark\ Exact: $243^{-0.6} = 3^{-3} = 1/27$ and $32^{-0.6} = 2^{-3} = 1/8$.

\textbf{Southern workers earn $2.25\times$ Northern workers; Northern capital earns
$3.4\times$ Southern capital.} That is the whole economics of the problem: each factor is
better paid where it is scarcer.

\begin{tcolorbox}[databox,title={Both interest rates are negative --- say so}]
$r_S = -6.52\%$ and $r_N = -3.00\%$. This is a property of the numbers the problem chose, not
an error: $\delta = 0.08$ exceeds the marginal product of capital in \emph{both} economies.
Economically it means both are over-capitalised --- they hold more capital than is worth
maintaining, and are dynamically inefficient in the sense of \S4.3. Neither would be in a
steady state with a sensible saving rate. Flag this rather than quietly reporting a negative
number.
\end{tcolorbox}

\item \textbf{Unification.} Capital and labour pool:
$$k_U = \frac{K_S + K_N}{L_S + L_N} = \frac{2430 + 160}{10 + 5} = \frac{2590}{15}
= \boxed{172.67}$$
$$y_U = 172.67^{0.4} = \boxed{7.850}$$
\checkmark

Note $k_U$ is the \textbf{labour-weighted average} of $k_S$ and $k_N$:
$\frac{10(243) + 5(32)}{15} = 172.67$. \checkmark

\textbf{Total output rises.} $Y_U = 7.850 \times 15 = 117.75$, against $90 + 20 = 110$
before --- a gain of $7.75$, about $7\%$, from nothing but reallocation.

\emph{Why is this not a violation of constant returns to scale?} CRS says that scaling
\emph{both} inputs by $\lambda$ scales output by $\lambda$. Here we are not scaling: we are
\textbf{pooling two economies with different capital--labour ratios}. Since $f$ is strictly
concave (Assumption 4.3), Jensen's inequality gives
$$f\!\left(\frac{L_S k_S + L_N k_N}{L_S+L_N}\right) > \frac{L_S f(k_S) + L_N f(k_N)}{L_S+L_N}$$
$$7.850 > \frac{90+20}{15} = 7.333 \quad\checkmark$$
\textbf{Equalising $k$ across a given total of capital and labour always raises total
output}, and the gain is larger the more unequal the initial ratios. This is the general case
for factor mobility.

\item \textbf{Factor prices in the unified country, and who wins.}
$$w_U = 0.6 \times 7.850 = \boxed{4.710}, \qquad
r^K_U = 0.4 \times 172.67^{-0.6} = \boxed{0.01819}, \qquad r_U = \boxed{-0.0618}$$
\checkmark

\begin{center}
\begin{tabular}{lrrl}
\toprule
& \textbf{Before} & \textbf{After} & \textbf{Effect} \\
\midrule
Southern workers      & $5.400$ & $4.710$ & \textbf{lose} $-12.8\%$ \\
Northern workers      & $2.400$ & $4.710$ & \textbf{gain} $+96.3\%$ \\
Southern capital ($r^K$) & $0.01481$ & $0.01819$ & \textbf{gain} $+22.8\%$ \\
Northern capital ($r^K$) & $0.05000$ & $0.01819$ & \textbf{lose} $-63.6\%$ \\
\bottomrule
\end{tabular}
\end{center}

\textbf{Interpretation.} Unification is a Kaldor--Hicks improvement --- total output rises by
$7\%$ --- but it is emphatically \textbf{not} a Pareto improvement. There are clear winners
and losers, and they are split \emph{by factor, not by country}:
\begin{itemize}[leftmargin=1.3em,itemsep=2pt,topsep=2pt]
\item \textbf{Winners:} Northern \emph{workers} (their wages nearly double) and Southern
      \emph{capital owners}.
\item \textbf{Losers:} Southern \emph{workers} and Northern \emph{capital owners}.
\end{itemize}
This is the standard political economy of economic integration and of immigration: the
scarce factor in each region loses when integration makes it less scarce. It explains why
South Korean workers might rationally oppose a unification that raises aggregate Korean
output --- and why the aggregate gain, being real, could in principle be used to compensate
them.

\item \textbf{Direction of movement: capital flows South $\to$ North, labour flows North
$\to$ South.}

\begin{itemize}[leftmargin=1.3em,itemsep=3pt,topsep=2pt]
\item \textbf{Machines move North.} Capital chases the higher return: $r^K_N = 0.050 >
      r^K_S = 0.0148$. An investor earns $3.4\times$ more per unit of capital in the North.
\item \textbf{People move South.} Labour chases the higher wage: $w_S = 5.4 > w_N = 2.4$.
\end{itemize}

Both flows do the same thing --- they \textbf{equalise $k$} --- and both stop exactly when
$k_S = k_N = 172.67$, at which point $w$ and $r^K$ are equal everywhere and there is no
further incentive to move. The model is \textbf{indifferent between the two channels}: only
$k$ matters, so whether adjustment happens through migration, through capital flows, or
through any mix of the two, the final allocation and all factor prices are identical.

\emph{What the model leaves out.} In reality the two channels are not interchangeable.
Capital moves faster and more cheaply than people; migration has social and political costs
the model does not price; and North Korean workers would not in fact be equally productive
with Southern capital, because human capital and institutions differ --- which is exactly the
gap Chapter 5 fills with \textbf{TFP differences}. The German reunification experience, where
massive capital transfers to the East did \emph{not} equalise productivity, is the obvious
cautionary case.

\end{enumerate}

\begin{tcolorbox}[theorybox,title={Verification code (executed)}]
\footnotesize
\begin{verbatim}
al, de = 0.4, 0.08
f  = lambda k: k**al
w  = lambda k: (1-al)*k**al
rK = lambda k: al*k**(al-1)

for name, K, L in [("South",2430,10), ("North",160,5), ("Unified",2590,15)]:
    k = K/L
    print(f"{name:8s} k={k:9.4f} y={f(k):7.4f} Y={f(k)*L:9.4f} "
          f"w={w(k):7.4f} rK={rK(k):8.5f} r={rK(k)-de:+.5f}")
# South    k= 243.0000 y= 9.0000 Y=  90.0000 w= 5.4000 rK= 0.01481 r=-0.06519
# North    k=  32.0000 y= 4.0000 Y=  20.0000 w= 2.4000 rK= 0.05000 r=-0.03000
# Unified  k= 172.6667 y= 7.8502 Y= 117.7537 w= 4.7101 rK= 0.01819 r=-0.06181
# gain from pooling: 117.7537 - 110 = +7.7537   (Jensen: f is strictly concave)
\end{verbatim}
\end{tcolorbox}
\end{tcolorbox}

%---- Topic: Alternative Production Technologies ----
\subsection{Alternative Production Technologies}

%------- 4.4 -------
% Topic: Alternative Production Technologies
% Keywords: AK model, endogenous growth, Inada conditions, constant marginal product, saving rate and growth
\begin{tcolorbox}[evenbox,title={Problem 4.4 \quad An AK Model}]
\begin{tcolorbox}[questionbox]
Suppose the production function takes the form $F(K,L) = AK$.\\[3pt]
(a) Which of the assumptions that we made about $F$ does this satisfy and which does it not
satisfy?\\
(b) Suppose every other assumption we used in the Solow model holds, and there is no
technological progress. Find an expression for $\dfrac{k_{t+1}}{k_t}$ (it may be useful to
follow the steps that lead to expression (4.2.2)).\\
(c) Will this economy grow in the long run? Explain.
\end{tcolorbox}

\textbf{Solution:}

\begin{enumerate}

\item \textbf{Checking the assumptions of \S4.1 one by one.}

\begin{center}
\begin{tabular}{llp{0.42\textwidth}}
\toprule
\textbf{Assumption} & \textbf{Status} & \textbf{Why} \\
\midrule
\textbf{4.1} CRS & \textbf{satisfied} & $F(\lambda K,\lambda L)=A\lambda K=\lambda AK=\lambda F(K,L)$ \\[3pt]
\textbf{4.2} Positive MP & \textbf{fails} & $F_K = A > 0$ \checkmark, but $F_L = 0$. Labour has \emph{zero} marginal product --- it is not a factor of production at all \\[3pt]
\textbf{4.3} Diminishing MP & \textbf{fails} & $F_{KK} = 0$, not $<0$. The marginal product of capital is \textbf{constant} at $A$ forever. ($F_{LL}=0$ too.) \\[3pt]
\textbf{4.4} Inada & \textbf{fails, both} & $\lim_{K\to0}F_K = A \ne \infty$ and $\lim_{K\to\infty}F_K = A \ne 0$ \\
\bottomrule
\end{tabular}
\end{center}

\textbf{Only CRS survives}, and Assumption 4.5 (exogenous population growth) is untouched
since it says nothing about $F$. The failure of \textbf{4.3} is the one that drives
everything below --- the absence of diminishing returns to capital is the entire content of
the AK model.

\item \textbf{Deriving $k_{t+1}/k_t$.} Follow the standard steps. Per worker,
$$y_t = \frac{Y_t}{L_t} = \frac{AK_t}{L_t} = A k_t$$
so $f(k) = Ak$ --- a straight line through the origin rather than a concave curve. The
capital accumulation equation is unchanged:
$$K_{t+1} = (1-\delta)K_t + sY_t$$
Divide by $L_{t+1} = (1+n)L_t$:
$$k_{t+1} = \frac{(1-\delta)k_t + s f(k_t)}{1+n} = \frac{(1-\delta)k_t + sAk_t}{1+n}$$
The $k_t$ factors out:
$$\boxed{\;\frac{k_{t+1}}{k_t} = \frac{1-\delta+sA}{1+n}\;}$$

\textbf{This is the key result: the growth factor is a constant, completely independent of
$k_t$.} Contrast with the standard Solow model, where $\frac{k_{t+1}}{k_t}
= \frac{(1-\delta) + s f(k_t)/k_t}{1+n}$ and the term $f(k)/k$ is \emph{decreasing} in $k$ ---
that decline is what produces convergence to a steady state. Here $f(k)/k = A$ is constant, so
the mechanism that generates a steady state is gone.

\item \textbf{Yes, this economy grows forever --- provided $sA$ is large enough --- and this
is \emph{endogenous} growth.}

The economy grows if and only if $\frac{k_{t+1}}{k_t} > 1$:
$$\frac{1-\delta+sA}{1+n} > 1 \iff 1-\delta+sA > 1+n
\iff \boxed{sA > n + \delta}$$
If this holds, $k$ (and hence $y = Ak$) grows at the \textbf{constant rate}
$$g_k = \frac{1-\delta+sA}{1+n} - 1 = \frac{sA - n - \delta}{1+n}$$
forever, with \textbf{no steady state and no convergence}. If $sA < n+\delta$ the economy
shrinks to zero instead. There is no interior rest point either way.

\textbf{Why the Inada conditions mattered.} In standard Solow, $\lim_{K\to\infty}F_K = 0$
guarantees that accumulation eventually runs into diminishing returns: each new machine adds
less, until investment only just covers depreciation and population growth. \textbf{Remove
diminishing returns and accumulation never runs out of steam.}

\textbf{The economically important consequence: the saving rate affects the long-run growth
rate, not just the level.}
$$\frac{\partial g_k}{\partial s} = \frac{A}{1+n} > 0$$
This is the sharpest possible contrast with Solow, where $s$ affects only $y\stst$ and has
\emph{no} effect on long-run growth (the ``level effect versus rate effect'' distinction from
\S4.2). A permanent increase in the saving rate here permanently raises growth. That makes
policy --- anything that raises $s$ or $A$ --- permanently consequential, which is why the AK
structure became the workhorse of the first generation of endogenous growth theory.

\textbf{Is it plausible?} The obvious objection is that $F_L = 0$ is absurd: labour clearly
matters. The standard defence is that $K$ should be read broadly, as \emph{physical plus
human capital plus knowledge}, and that the constant marginal product proxies for
\textbf{externalities} --- one firm's capital raises other firms' productivity, so the
\emph{social} return to capital does not diminish even though the private return does. The
model is best read as a reduced form for that idea rather than as a literal technology.

\end{enumerate}

\begin{tcolorbox}[theorybox,title={Verification code (executed)}]
\footnotesize
\begin{verbatim}
A, s, delta, n = 0.5, 0.3, 0.08, 0.01
ratio = (1 - delta + s*A) / (1 + n)          # 1.05940594 -- constant

k = 1.0
for t in range(200):
    k = ((1-delta)*k + s*A*k) / (1+n)
print(k, 1.0*ratio**200)   # 102915.498144  102915.498144  -> exactly geometric

print(s*A, delta + n + n*delta)              # 0.15 > 0.0908  -> grows
print(((1-delta+0.4*A)/(1+n)) - 1)           # s: 0.3->0.4 raises growth 5.94% -> 10.89%
\end{verbatim}
\end{tcolorbox}
\end{tcolorbox}

%------- 4.5 -------
% Topic: Alternative Production Technologies
% Keywords: Malthusian model, endogenous population, subsistence, diminishing returns to labour, poverty trap
\begin{tcolorbox}[oddbox,title={Problem 4.5 \quad A Malthusian Model}]
\begin{tcolorbox}[questionbox]
Look up Thomas Malthus and read a little about his work. We are going to formalize a simple
version of his ideas in a little model. The production function is
$$Y_t = A_t L_t^{\alpha} \qquad\text{(4.5.5)}$$
where $A_t$ is technology, $L$ is labor supply and $\alpha < 1$. There is no capital, no
saving and no investment, so consumption is given by $C_t = Y_t$ (4.5.6). The population
evolves according to
$$L_{t+1} = L_t\left(1 + \gamma\left(\frac{C_t}{L_t} - c\right)\right) \qquad\text{(4.5.7)}$$
where $\gamma$ and $c$ are parameters.\\[3pt]
(a) Malthus had in mind an agricultural economy where the total amount of land is fixed. What
does that have to do with the way we wrote down the production function? What is the
significance of assuming $\alpha < 1$?\\
(b) What does equation (4.5.7) mean? What is $c$? What is $\gamma$? Why did Malthus think
that something like equation (4.5.7) applied?\\
(c) Assume $A_t = A$ is constant over time. Using (4.5.5), (4.5.6) and (4.5.7), find an
expression for the percentage rate of population growth $\frac{L_{t+1}}{L_t}$ as a function of
the level of population $L_t$. Plot this function.\\
(d) Assume $A_t = A$ is constant over time. Will GDP per capita grow in the long run? Will the
population grow in the long run?\\
(e) Suppose there is a one-time increase in $A_t$, from $A$ to $A' > A$. What will happen to
GDP per capita and the size of the population in the short run and in the long run? Show your
reasoning graphically and/or algebraically.\\
(f) Suppose now that instead of being constant, technology improves at a constant rate, so
that $A_{t+1} = A_t(1+g)$. Will GDP per capita grow in the long run? Will the population grow
in the long run? Explain.\\
(g) When we have a constant rate of technological progress, how does the long-run level of
consumption per capita depend on $\gamma$? Why?
\end{tcolorbox}

\textbf{Solution:}

\begin{enumerate}

\item \textbf{Land is the hidden fixed factor, and $\alpha<1$ is what encodes it.}

Write the ``true'' technology with land $T$ explicitly and constant returns in
(labour, land): $Y = \tilde A\, L^{\alpha} T^{1-\alpha}$. If land is \textbf{fixed} at
$\bar T$, absorb it into the constant:
$$Y = \underbrace{\tilde A \bar T^{1-\alpha}}_{\equiv\,A} L^{\alpha}$$
which is exactly (4.5.5). So the reduced-form production function \emph{is} the
constant-returns technology with the fixed factor substituted out.

\textbf{Significance of $\alpha < 1$:} it delivers \textbf{diminishing returns to labour}.
$$\pd{Y}{L} = \alpha A L^{\alpha-1} > 0, \qquad
\frac{\partial^2 Y}{\partial L^2} = \alpha(\alpha-1)AL^{\alpha-2} < 0$$
and, critically, output per worker
$$\frac{Y_t}{L_t} = A L_t^{\alpha-1}$$
is \textbf{decreasing} in $L$ because $\alpha - 1 < 0$. More people working the same fixed
land means less output each. \textbf{This is the engine of the entire model} --- without it
($\alpha = 1$) output per head would be independent of population and nothing that follows
would work.

\item \textbf{Equation (4.5.7) makes population growth endogenous, responding to living
standards.}

Rewrite it as
$$\frac{L_{t+1} - L_t}{L_t} = \gamma\left(\frac{C_t}{L_t} - c\right)$$
so the population \emph{growth rate} is proportional to the excess of consumption per capita
over $c$.

\begin{itemize}[leftmargin=1.3em,itemsep=3pt,topsep=2pt]
\item \textbf{$c$ is the subsistence level of consumption.} If $C/L > c$, the population
      grows; if $C/L < c$, it shrinks; if $C/L = c$, it is constant. It is the living
      standard at which the population exactly reproduces itself.
\item \textbf{$\gamma$ measures the responsiveness of population growth to living standards}
      --- the semi-elasticity of demographic behaviour with respect to income. A high $\gamma$
      means population reacts strongly and quickly; a low $\gamma$ means it is sluggish.
\end{itemize}

\textbf{Why Malthus believed it.} Writing in 1798, he had in mind two mechanisms:
\begin{itemize}[leftmargin=1.3em,itemsep=2pt,topsep=2pt]
\item \textbf{Positive checks} operating below subsistence --- famine, disease and war raise
      the death rate, especially infant and child mortality.
\item \textbf{Preventive checks} operating above it --- prosperity brings earlier marriage
      and higher fertility; hardship delays marriage and reduces it.
\end{itemize}
This was an accurate description of pre-industrial demography, where fertility genuinely did
respond to harvests and real wages. Note the sharp contrast with Kurlat's \textbf{Assumption
4.5}, which makes $n$ exogenous --- the book flags in \S4.1 that this is ``actually a very
big deal'', and this exercise is where you see why.

\item \textbf{The population growth function.} Substitute (4.5.5) and (4.5.6):
$$\frac{C_t}{L_t} = \frac{Y_t}{L_t} = \frac{A L_t^{\alpha}}{L_t} = A L_t^{\alpha-1}$$
Therefore
$$\boxed{\;\frac{L_{t+1}}{L_t} = 1 + \gamma\left(A L_t^{\alpha-1} - c\right)\;}$$

\textbf{Shape of the plot} ($L$ on the horizontal axis, $L_{t+1}/L_t$ on the vertical):
\begin{itemize}[leftmargin=1.3em,itemsep=2pt,topsep=2pt]
\item \textbf{Strictly decreasing} in $L$, since $\alpha - 1 < 0$.
\item As $L \to 0^+$: $L^{\alpha-1}\to\infty$, so the growth factor $\to +\infty$.
\item As $L \to \infty$: $L^{\alpha-1}\to 0$, so the growth factor $\to 1 - \gamma c$
      (a horizontal asymptote below 1 --- population shrinks).
\item It crosses the horizontal line at height $1$ exactly once, at $L\stst$ below.
\end{itemize}
Draw the horizontal line at $1$: the crossing point is the steady state, growth is above it
to the left and below it to the right, which shows the dynamics are \textbf{globally stable}.

\item \textbf{Neither GDP per capita nor population grows in the long run.}

Set $\frac{L_{t+1}}{L_t} = 1$:
$$\gamma\left(A(L\stst)^{\alpha-1} - c\right) = 0
\;\Longrightarrow\; A (L\stst)^{\alpha-1} = c
\;\Longrightarrow\; \boxed{L\stst = \left(\frac{A}{c}\right)^{\frac{1}{1-\alpha}}}$$

\textbf{Stability.} For $L < L\stst$, $AL^{\alpha-1} > c$, so $L$ grows; for $L > L\stst$ it
shrinks. Globally stable, converging monotonically. \checkmark\ (Simulation from $L_0 = 500$
with $A=100$, $c=2$, $\alpha=0.5$ converges to $L\stst = 2500$ exactly.)

\textbf{And here is the punchline:} at the steady state
$$\frac{C}{L} = A (L\stst)^{\alpha-1} = c$$
\textbf{GDP per capita equals subsistence, always.} It cannot be anything else --- if it were
above, population would grow and push it down; if below, population would fall and push it
up. So:
\begin{itemize}[leftmargin=1.3em,itemsep=2pt,topsep=2pt]
\item \textbf{GDP per capita: no long-run growth.} It is pinned at $c$.
\item \textbf{Population: no long-run growth.} It converges to $L\stst$.
\end{itemize}

\item \textbf{A one-time increase $A \to A'$: living standards jump, then fall all the way
back. Only population is permanently higher.}

\begin{itemize}[leftmargin=1.3em,itemsep=4pt,topsep=2pt]
\item \textbf{Short run.} $L$ is predetermined and cannot jump. So consumption per capita
      jumps immediately:
      $$\frac{C}{L} = A' (L\stst)^{\alpha-1} > c$$
      \checkmark\ In simulation, doubling $A$ from 100 to 200 makes $C/L$ jump from $2.00$ to
      $4.00$ on impact. \textbf{The generation alive at the moment of the innovation is
      genuinely better off.}
\item \textbf{Transition.} Now $C/L > c$, so by (4.5.7) the population grows. As $L$ rises,
      diminishing returns push $C/L = A'L^{\alpha-1}$ back down.
\item \textbf{Long run.} The process stops only when $C/L$ is back at $c$:
      $$L\stst_{\text{new}} = \left(\frac{A'}{c}\right)^{\frac{1}{1-\alpha}}
      = L\stst \left(\frac{A'}{A}\right)^{\frac{1}{1-\alpha}} > L\stst$$
      $$\frac{C}{L}\Big|_{\text{long run}} = c \quad \text{--- exactly as before.}$$
      \checkmark\ Simulation: $A$ doubles, $L\stst$ goes $2500 \to 10{,}000$ (a factor of
      $4 = 2^{1/(1-0.5)}$, as the formula requires), and $C/L$ returns to $2.000000$.
\end{itemize}

\textbf{This is the Malthusian trap.} Technological progress is converted entirely into
\emph{more people}, not into better-off people. It is the single most important idea in the
exercise, and it explains the central fact of Chapter 3: living standards were roughly flat
for millennia despite continuous, if slow, technological improvement. Better ploughs, crop
rotation and the horse collar all raised $A$ --- and all produced a larger population living
at the same subsistence level.

\item \textbf{With sustained progress $A_{t+1} = A_t(1+g)$: population grows forever, GDP per
capita still does not.}

On a balanced path $C/L$ must be constant. Since $C_t/L_t = A_t L_t^{\alpha-1}$, constancy
requires $A_t L_t^{\alpha-1}$ constant, i.e.\ $L_t^{1-\alpha} \propto A_t$, so
$$L_t \propto A_t^{\frac{1}{1-\alpha}}
\quad\Longrightarrow\quad
\boxed{\;\frac{L_{t+1}}{L_t} = (1+g)^{\frac{1}{1-\alpha}}\;}$$
\checkmark\ Simulated growth rate with $g = 0.02$, $\alpha = 0.5$: $0.04040000$, against the
theoretical $(1.02)^{2}-1 = 0.04040000$. Exact match.

\textbf{Population growth is \emph{faster} than technological progress} --- amplified by
$\frac{1}{1-\alpha} > 1$, because it takes a more-than-proportional increase in population to
absorb a given productivity gain when returns diminish.

\textbf{And GDP per capita?} Constant --- it does not grow. But note it settles \emph{above}
subsistence. Setting $L_{t+1}/L_t = (1+g)^{1/(1-\alpha)}$ in (4.5.7):
$$(1+g)^{\frac{1}{1-\alpha}} = 1 + \gamma\left(\frac{C}{L} - c\right)
\;\Longrightarrow\;
\boxed{\;\frac{C}{L} = c + \frac{(1+g)^{\frac{1}{1-\alpha}} - 1}{\gamma}\;}$$
\checkmark\ Simulated: $2.808000$; formula with $c=2$, $\gamma=0.05$: $2 + 0.0404/0.05
= 2.808000$. Exact match.

So sustained technological progress buys a permanently \emph{higher} standard of living than
the no-growth case --- but a \textbf{constant} one. The economy escapes bare subsistence but
never achieves growth in income per head. Only breaking the link between living standards and
fertility (i.e.\ $\gamma \to 0$, the demographic transition) allows modern growth.

\item \textbf{Higher $\gamma$ means \emph{lower} long-run consumption per capita.}

From the expression just derived,
$$\pd{}{\gamma}\left(\frac{C}{L}\right)
= -\frac{(1+g)^{\frac{1}{1-\alpha}} - 1}{\gamma^{2}} < 0$$

\checkmark\ Verified by simulation across a range of $\gamma$ (with $g=0.02$, $\alpha=0.5$,
$c=2$):

\begin{center}
\begin{tabular}{crr}
\toprule
$\gamma$ & \textbf{Formula} & \textbf{Simulated} \\
\midrule
$0.02$ & $4.020000$ & $4.020000$ \\
$0.05$ & $2.808000$ & $2.808000$ \\
$0.10$ & $2.404000$ & $2.404000$ \\
$0.25$ & $2.161600$ & $2.161600$ \\
$0.50$ & $2.080800$ & $2.080800$ \\
\bottomrule
\end{tabular}
\end{center}

\textbf{Why.} The logic runs backwards from the required population growth rate.
Technological progress at rate $g$ \emph{forces} the population to grow at
$(1+g)^{1/(1-\alpha)}$ on the balanced path --- that rate is determined by technology and
$\alpha$, and $\gamma$ has no say in it. But equation (4.5.7) says population growth is
\emph{produced} by living standards above subsistence. So living standards must sit at
whatever level generates the required growth:
\begin{itemize}[leftmargin=1.3em,itemsep=3pt,topsep=2pt]
\item If $\gamma$ is \textbf{high} (population very responsive), only a \emph{small} excess
      over subsistence is needed to generate that much population growth. Living standards
      settle close to $c$.
\item If $\gamma$ is \textbf{low} (population sluggish), a \emph{large} excess is needed.
      Living standards settle well above $c$.
\end{itemize}
As $\gamma \to 0$ the required excess becomes unbounded and the Malthusian mechanism
disappears entirely --- population no longer responds to income, and technological progress
finally shows up as growth in income per head instead of in numbers.

\textbf{The historical reading.} This is exactly the \textbf{demographic transition}. What
allowed Europe to escape the Malthusian trap after about 1800 was not only faster
technological progress but a \emph{collapse in $\gamma$}: fertility stopped rising with
income and began falling with it. In this model, that is the difference between the world of
Chapter 3's flat pre-1800 living standards and the world the Solow model of \S4.5 is designed
to describe.

\end{enumerate}

\begin{tcolorbox}[theorybox,title={Verification code (executed)}]
\footnotesize
\begin{verbatim}
import numpy as np
alpha, c_sub, gamma = 0.5, 2.0, 0.05

def simulate(A_series, L0, T):
    L, cpc = np.empty(T+1), np.empty(T+1)
    L[0] = L0
    for t in range(T+1):
        cpc[t] = A_series[t] * L[t]**(alpha-1)          # C/L = A L^(alpha-1)
        if t < T:
            L[t+1] = L[t] * (1 + gamma*(cpc[t] - c_sub))
    return L, cpc

# (d) A constant -> L* = (A/c)^(1/(1-alpha)) = 2500, C/L -> c
L, cpc = simulate(np.full(401, 100.0), L0=500.0, T=400)   # L[-1]=2500.0, cpc[-1]=2.0

# (e) one-time A: 100 -> 200 : C/L jumps 2.0 -> 4.0, returns to 2.0; L: 2500 -> 10000
# (f) A grows at g=2% : L grows at (1+g)**(1/(1-alpha))-1 = 0.0404 exactly
#     C/L -> c + ((1+g)**(1/(1-alpha)) - 1)/gamma = 2 + 0.0404/0.05 = 2.808
\end{verbatim}
\end{tcolorbox}
\end{tcolorbox}

%======================================================================
\section*{Appendix \quad Cross-reference to the course}
\addcontentsline{toc}{section}{Appendix --- Cross-reference to the course}
%======================================================================

\begin{tcolorbox}[theorybox,title={Which exercises map to Lista de Exerc\'icios 1}]
The Lista (issued 03/08/2026, due 10/08/2026, 10 points) draws directly on this material:

\begin{center}
\begin{tabular}{clll}
\toprule
\textbf{Q} & \textbf{Pts} & \textbf{Source} & \textbf{See} \\
\midrule
1 & 3 & Kurlat ch.\ 1, authored problem & Problem 1.2 (same structure, Brazil/Colombia) \\
2 & 2 & \textbf{Kurlat Exercise 2.2} & Problem 2.2 \\
3 & 2 & \textbf{Kurlat Exercise 3.3} & Problem 3.3 \\
4 & 3 & Kurlat ch.\ 4, extended & Problem 4.1 (plus the $n$ halving extension) \\
\bottomrule
\end{tabular}
\end{center}

Questions 2 and 3 are the book's own exercises assigned \emph{by number}. Question 1 mirrors
Problem 1.2 with Brazilian and Colombian data instead of US and Thai. Question 4 is
Problem 4.1 with an added part.

\textbf{Note that Q3 is a data exercise} --- it requires downloading the Penn World Table and
producing three regional scatterplots. Do not leave it to the last day.
\end{tcolorbox}

\begin{tcolorbox}[theorybox,title={Course rule files}]
\begin{center}
\begin{tabular}{lll}
\toprule
\textbf{Chapter} & \textbf{Aula} & \textbf{Rules file} \\
\midrule
1--2 & 1 & \texttt{rules/01\_mensuracao\_agregados.md} \\
3, 4.1--4.2 & 2 & \texttt{rules/02\_crescimento\_solow.md} \\
4.3--4.5 & 3 & \texttt{rules/03\_solow\_evidencias.md} \\
\bottomrule
\end{tabular}
\end{center}
Page-level anchors for every section are in \texttt{Map/registro-livros.md}. Kurlat's printed
page numbers equal its PDF page numbers (offset $0$), so \texttt{/study-pack} can extract any
range directly.
\end{tcolorbox}

\end{document}
